% Sub-stage F4: the chunked first stage of the default build and the rows of the C code: the chunk
% RATPOINTS_CHUNK and both arms of the first stage, the rows of m + chunk - 1 bit arrays with their
% copies, the builders of init.c and the generated tables, the tables of a prime power as words, and
% the rows kept in the tables of the sieve entries, built before a denominator by the loop of sift.
% Proposed names are set with \lname and become \lean once the declarations exist.
\chapter{The chunked first stage and the rows}\label{chap:rows}

F4 models the first stage of the sieve as the default build of ratpoints compiles it, and the rows
of the moduli as ratpoints builds and keeps them. The C code has two arms of the first stage, which
the preprocessor chooses by the number \verb|RATPOINTS_CHUNK| of registers of a chunk: for a chunk
from $2$ to $16$, the chunked arm, which loads a chunk of bit arrays into registers, and-s the rows
of the moduli into them from a start that it steps through each row, and does the bit arrays after
the last whole chunk in legs of $8$, $4$, $2$ and $1$; for the other chunks, the arm that and-s one
row at a time into the bit arrays of a piece, which the model has had since stage~C2 and which is
the build \verb|chunk1| that ratpoints tests. The chunked arm reads up to $\mathrm{chunk} - 1$ bit
arrays past the first $p$ bit arrays of a row of the modulus $p$, so every row of ratpoints has
$p + \mathrm{chunk} - 1$ bit arrays, the last ones copies of the first. F4 makes the chunk a
constant of the model, $16$ as in the default build, models both arms and chooses between them as
the conditional compilation of the C code does; the proofs hold for every value of the chunk, and
the build with other values of the constants varies it.

Until now the model made each row from the admission test of its sieve entry, one period laid out in
words by one function for every kind of modulus, and kept it in a thunk, a suspended computation. F4
builds the rows as ratpoints does: the row of $0$ of a prime is the pattern \verb|sieves0| that the
build generates, the other rows of a prime come from the function that the build generates for the
prime from \verb|CODE_INIT_SIEVE1| or \verb|CODE_INIT_SIEVE2|, the rows of a prime power from
\verb|_ratpoints_sieve_init_power| with \verb|lay_out_pattern|, and those of a product from
\verb|_ratpoints_sieve_init_product|, out of the rows of its factors. They are kept where ratpoints
keeps them, in the tables of rows of the sieve entries: one table of rows for each entry of
ratpoints, reached as the C code reaches the entry (a prime by itself, a prime power by its prime
and its exponent, a product by its position among the moduli), the rows that a denominator needs
built before it by the loop of \verb|sift|, a product building the missing rows of its factors into
their tables of rows first, and the tables of rows handed on through the run as state. With them
come the tables of a prime power as bit arrays of words, as \verb|examine_power| fills them, and the
forbidden divisors of the primes of the table read their words from \verb|sieves0|, as
\verb|sieving_info| does. In this sub-stage, as in the others, the model follows the data and the
loops of the C code where the C code has them; the speed of the model is not an aim.

The rows that the builders build are the closed forms of the admission tests: their first $p$ bit
arrays are the ones the plain builder computed, and the $\mathrm{chunk} - 1$ after them are copies
of the first ones (F4.3 gives the plain row these copies too, until F4.10 retires it); the chunked
arm and-s the same words into a piece as the other arm, since its reads past the first $p$ bit
arrays of a row meet copies of the bit arrays that the other arm reads; and every row in the tables
of rows of the run is the closed form of the row of its entry, so that the rows that a denominator
reads are those that the thunks of the rows gave. So the output of \verb|findPoints| does not
change, nor do the statements that comparator certifies, \verb|Params| or \verb|Params.WF|: the
well-formedness already makes every prime of the moduli odd and every modulus a product of primes of
the table below $\mathrm{maxPrimeEven}$, which is all the builders need. The contracts of the stages
keep their conclusions; a sieve spec that fits also has the end $p$ of the walk through its row and
a row with its copies; the contracts of one denominator, of the run and of \verb|sieving_info| take,
or give, a store of the tables of rows that fits the sieve entries, in place of the thunks of the
rows. Comparator checks the statements again, since the proofs change.

\section{What the C code does}\label{sec:rowsC}

\paragraph{The chunk.} Each build defines \verb|RATPOINTS_CHUNK| unless the compiler is given it:
$16$ for the builds with vector registers, the default build (\verb|USE_AVX|) among them
(\csource{rp-private.h}{160}{165}, \csource{rp-private.h}{222}{224},
\csource{rp-private.h}{248}{250}, \csource{rp-private.h}{274}{276}), and $1$ for the build with
$64$-bit bit arrays (\csource{rp-private.h}{293}{295}). The largest array size is the largest
multiple of the chunk up to $2^{20}$ (\csource{find_points.c}{124}{124}), and the set-up
\verb|find_points_work_1| rounds the array size, after its two caps, up to a multiple of the chunk
(\csource{find_points.c}{4666}{4677}), so that of the pieces of an interval only the last one can
end in a partial chunk.

\paragraph{The chunked first stage.} For a chunk from $2$ to $16$ (\csource{sift.c}{575}{575}),
\verb|_ratpoints_sift0| sets for each of the first $s_1$ moduli $p$, those of the first stage, the
start $\mathrm{ptr} + ((w_{\mathrm{low}} + \mathrm{offset}) \bmod p)$ of its walk through its row,
by \verb|mod_mul| under \verb|small| and by \verb|mod| otherwise (\csource{sift.c}{596}{606}); with
no modulus of the first stage, as \verb|-n 0| asks, it writes the pattern of the class,
\verb|bits64|, whose words are all the class word, into every bit array of the piece, from the last
(\csource{sift.c}{608}{615}). Otherwise it takes the bit arrays in chunks of \verb|RATPOINTS_CHUNK|
up to the end $w_{\mathrm{high,full}}$ of the whole chunks (\csource{sift.c}{616}{624}): the first
modulus loads the registers \verb|reg0|, \dots, each the class word and-ed with a bit array of its
row from the start on, and moves its start on by the chunk and back by $p$ while it is at least the
end $\mathrm{ptr} + p$ (\csource{sift.c}{633}{684}); each further modulus and-s its bit arrays into
the registers and moves its start in the same way (\csource{sift.c}{744}{806}); and the registers
are stored into the bit arrays of the chunk (\csource{sift.c}{864}{912}). The code for each register
is written out under \verb|#if (RATPOINTS_CHUNK >= k)|, for $k$ up to $16$, with output under
\verb|DEBUG| after each modulus. The $t < \mathrm{chunk}$ bit arrays after the last whole chunk go
in legs of $8$, $4$, $2$ and $1$ registers, one for each set bit of $t$, the widest first, each from
the part of $t$ above its own bit, $0$, $t \mathbin{\&} 8$, $t \mathbin{\&} 12$ and $t \mathbin{\&}
14$, the first three only for a chunk above $8$, $4$ and $2$ (\csource{sift.c}{915}{940}); a leg,
\verb|tail_leg| (\csource{sift.c}{473}{495}), loads at most eight registers from the first modulus,
and-s the others in and stores them, reading the rows past the starts and leaving the starts as they
are. A start is below $p$ after each chunk, so a chunk or a tail reads up to bit array $p +
\mathrm{chunk} - 2$ of a row, one of the copies at its end. The check that the chunk is at most
$16$, an \verb|#error| (\csource{sift.c}{927}{929}), sits inside the conditional compilation of the
chunked arm and never fires: a chunk above $16$ compiles the other arm
(\csource{sift.c}{943}{1022}), which writes \verb|bits64| into the bit arrays from the last
(\csource{sift.c}{954}{954}) and and-s one row at a time.

\paragraph{The rows of a prime.} The row of a residue $b$ modulo the prime $p$ has $p +
\mathrm{chunk} - 1$ bit arrays, bit $t$ of word $j$ standing for the bit number $64 j + t$. Its row
of $0$ is the row of $p$ in the table \verb|sieves0| that the build generates
(\csource{gen_find_points_h.c}{102}{134}; included at \csource{find_points.c}{37}{37}): the $p$
words with the bits of the multiples of $p$ cleared, written $\mathrm{rbaPack}$ times one after the
other, so that word $j$ is word $j \bmod p$, and the copies. For $b \ne 0$ the build generates one
function for each prime (\csource{gen_init_sieve_h.c}{48}{62}), from \verb|CODE_INIT_SIEVE1| for $p
< 64$ (\csource{init.c}{118}{183}) and from \verb|CODE_INIT_SIEVE2| for $p > 64$
(\csource{init.c}{186}{244}), with its table \verb|sieve_init| of pointers to them
(\csource{find_points.c}{133}{133}). Both compute the pattern of $p$ bits, bit $a$ the entry of
\verb|is_f_square| at $a\, b^{-1} \bmod p$, four bits at a time with four residues stepped by $4
b^{-1} \bmod p$ (\csource{init.c}{42}{46}, \csource{init.c}{62}{83}) and a plain loop for the rest.
\verb|CODE_INIT_SIEVE1| repeats the pattern through a word by or-ing it in shifted by $p$, $2p$,
$4p$, \dots, and rotates the first $p$ words of the row out of two words; \verb|CODE_INIT_SIEVE2|
keeps the pattern in $\lfloor p/64 \rfloor + 1$ words and rotates the row out of a ring of them, by
$p \bmod 64$, which is not $0$ for an odd prime. Then \verb|RP_INIT_REPLICATE| copies the first $p$
words into the next $(\mathrm{rbaPack} - 1)\, p$, so that word $j$ is word $j \bmod p$, and the
first $(\mathrm{chunk} - 1)\,\mathrm{rbaPack}$ words to the end (\csource{init.c}{99}{111}), and the
function stores the row into the entry with \verb|RP_PUBLISH| and returns it.

\paragraph{The rows of a prime power and of a product.} \verb|_ratpoints_sieve_init_power|
(\csource{init.c}{359}{393}) clears a pattern of \verb|RP_MODWORDS| words and sets bit $x$ when bit
$x\, b^{-1}$ of \verb|fsq| is set, for a unit $b$, and bit $x$ of a unit $x$ when bit $b\, x^{-1}$
of \verb|gsq| is set, for $p \mid b$; \verb|lay_out_pattern| (\csource{init.c}{259}{297}) repeats
the pattern of period $m$ through enough words and cuts word $x$ of the row out of them at the
offset $64 x \bmod m$, stepped by $64 \bmod m$, then appends the copies.
\verb|_ratpoints_sieve_init_product| (\csource{init.c}{399}{445}) reads the row of $b \bmod q$ of
each factor of modulus $q$ from the entry of the factor, building it first by the function of that
entry when it is not there (\csource{init.c}{407}{414}); then it lays the row of the first factor
down in blocks of $q$ bit arrays, and-s in those of the others in blocks of theirs, appends the
copies, and stores the row.

\paragraph{Where the rows live.} A sieve entry has the array \verb|sieve| of the rows of its
residues and the pointer \verb|init| to the function that builds them
(\csource{rp-private.h}{542}{555}; the type of \verb|init| is \verb|ratpoints_init_fun|,
\csource{rp-private.h}{497}{498}). \verb|examine_prime| makes the entry of a prime that says
something with its function from \verb|sieve_init| (\csource{find_points.c}{1800}{1800}), its row of
$0$ from \verb|sieves0| and the other rows absent (\csource{find_points.c}{1820}{1821});
\verb|make_power| makes the entry of a prime power once, when a modulus first needs it, with the
builder of the rows of a prime power and all rows absent (\csource{find_points.c}{2414}{2429});
\verb|make_modulus| takes it for a prime power and makes a product its own entry, with the builder
of the rows of a product, all rows absent, and pointers to the entries of its factors, those of the
primes from \verb|examine_prime| and those of the prime powers from \verb|make_power|
(\csource{find_points.c}{2434}{2475}). So a prime or a prime power has one table of rows, shared by
the modulus and by the products that have it as a factor; ratpoints never meets that sharing within
a run, since its moduli of the first two stages share no prime (\verb|take_entries| drops a
candidate whose primes meet those of the moduli taken, both calls with one mask,
\csource{find_points.c}{2490}{2520}; \verb|rp-private.h| says so at lines 513-514), so that each
table of rows is reached through one modulus only, the prime or the prime power itself or one
product that has it as a factor, whereas well-formed parameters of the model may reach one table of
rows through several. \verb|sift| sets up a denominator by \verb|set_rows|
(\csource{find_points.c}{3149}{3170}): for each modulus of the first two stages the row of its
residue from the table of rows of its entry, perhaps absent, and the end $\mathrm{ptr} + p$, and
whether a row is absent; then for each absent row it reads the pointer again under the lock of the
tables, builds the row by the function of the entry when it is still absent, and sets the pointer
and the end (\csource{find_points.c}{3522}{3554}). The rows stay in the tables of rows for the rest
of the run.

\paragraph{The tables of a prime power.} \verb|examine_power| keeps the squares modulo $m = p^e$ and
the tables \verb|fsq| and \verb|gsq| as bit arrays of \verb|RP_MODWORDS| words
(\csource{rp-private.h}{523}{523}, \csource{rp-private.h}{531}{534}), clears them in one loop and
sets their bits one at a time (\csource{find_points.c}{2282}{2300}).

\paragraph{Why this changes nothing.} Write $\mathrm{closedRow}(\mathrm{adm}, m, r)$ for the array
of $(m + \mathrm{chunk} - 1)\,\mathrm{rbaPack}$ words whose word $j$ has bit $t$ set exactly when
$\mathrm{adm}(r)$ admits $64 j + t$. Each builder builds the closed form of the admission test of
its entry at its residue, for a prime at a residue other than $0$, the row of $0$, \verb|sieves0|,
being the closed form for every table of the prime; the test is periodic with period $m$, so that
the copies at the end of a row are its first bit arrays. The chunked arm reads, for the bit array $c
+ j$ of a piece and a modulus with the offset $o$, bit array $s + j$ of the row, where $s = (w_0 + o
+ c) \bmod p$ and $j < \mathrm{chunk}$; that is a copy of bit array $(w_0 + o + c + j) \bmod p$,
which the other arm reads; and the legs of a tail of $t < 16$ bit arrays cover the bit arrays $0,
\dots, t - 1$ after the whole chunks once each. A row of the tables of rows is either one that the
run put there, the row of $0$ of a prime or a row that a builder built, or absent; so every row
present is the closed form of the row of its entry, and two entries with one table of rows, a prime
or a prime power as a modulus and as a factor, have the same table of the prime or tables of the
prime power, and so the same closed forms. The bits of \verb|fsq| and \verb|gsq| are the Booleans
that they replace, and the first $p$ words of \verb|sieves0| are the words of the forbidden divisor
$p$. So F4 is a data refinement: the sieve computes the same predicate as before.

\section{The chunk}\label{sec:chunkConstants}

\begin{definition}[The chunk]\label{def:chunk}
  \lean{Ratpoints.chunk, Ratpoints.chunk_pos}
  \leanok
  The number of bit arrays that the chunked first stage takes at a time, $\mathrm{chunk} = 16$, the
  value of the default build (\verb|RATPOINTS_CHUNK|: \csource{rp-private.h}{160}{165},
  \csource{rp-private.h}{222}{224}, \csource{rp-private.h}{248}{250},
  \csource{rp-private.h}{274}{276}, \csource{rp-private.h}{293}{295}, all the definitions of the
  builds, as $\mathrm{rbaPack}$ owns those of its macro), with the fact $0 < \mathrm{chunk}$ checked
  by evaluation, which the rounding of the array size uses. No proof uses its value: the chunked
  first stage needs $\mathrm{chunk} \le 16$, which the test that chooses the arm gives
  (Definition~\ref{def:siftPiece}). The largest array size becomes
  $\lfloor 2^{20}/\mathrm{chunk} \rfloor \cdot \mathrm{chunk}$ and the array size is rounded up to a
  multiple of the chunk after its caps (Definition~\ref{def:arraySize}); every row has
  $p + \mathrm{chunk} - 1$ bit arrays (Definition~\ref{def:closedRow}).
\end{definition}

\section{The closed form of a row}\label{sec:closedForm}

\begin{definition}[The closed form of a row]\label{def:closedRow}
  \uses{def:chunk, def:rbaPack, def:rows}
  \lean{Ratpoints.Sieve.closedRow, Ratpoints.Sieve.wordOfBits, Ratpoints.Sieve.SieveSpec.Wraps}
  \leanok
  For a test $\mathrm{adm}$ of a residue and a numerator, a modulus $m$ and a residue $r$: the array
  of $(m + \mathrm{chunk} - 1)\,\mathrm{rbaPack}$ words, $m + \mathrm{chunk} - 1$ bit arrays, whose
  word $j$ has bit $t$ set exactly when $t < 64$ and $\mathrm{adm}(r)$ admits $64 j + t$, each word
  built bit by bit from its $64$ bits. With the admission test of a sieve entry or of a factor
  (Definitions~\ref{def:admits}, \ref{def:examinePower} and~\ref{def:SieveEntry}) it is the row that
  ratpoints builds for the residue, against which the builders are proved, as the admission tests
  were for the plain builder. The row of a sieve spec \emph{wraps} when its bit arrays
  $p, \dots, p + \mathrm{chunk} - 2$ are copies of its first ones: word $\kappa < \mathrm{rbaPack}$
  of bit array $x < p + \mathrm{chunk} - 1$ is word $\kappa$ of bit array $x \bmod p$.
\end{definition}

\begin{lemma}[The closed form]\label{lem:closedRow}
  \uses{def:closedRow, def:SieveEntry, def:row}
  \lean{Ratpoints.Sieve.getLsbD_closedRow, Ratpoints.Sieve.eq_closedRow, Ratpoints.Sieve.size_closedRow, Ratpoints.Sieve.closedRow_congr, Ratpoints.Sieve.getElem!_closedRow_mod, Ratpoints.Sieve.wraps_closedRow, Ratpoints.Sieve.SieveEntry.wraps_closedRow, Ratpoints.Sieve.getLsbD_wordOfBits, Ratpoints.Sieve.PatternBits, Ratpoints.Sieve.RowWords, Ratpoints.Sieve.RowWords.toArray_eq_closedRow, Ratpoints.Sieve.SieveEntry.row_eq_closedRow}
  \leanok
  Bit $t$ of word $j < (m + \mathrm{chunk} - 1)\,\mathrm{rbaPack}$ of the closed form is $t < 64$
  and $\mathrm{adm}(r)(64 j + t)$, and an array of that size with these bits is the closed form. For
  a test of period $m > 0$ in the numerator, word $j$ of the closed form is word
  $j \bmod m\,\mathrm{rbaPack}$, so the closed form of the row of a sieve entry with a positive
  modulus wraps, whatever the offset, the start and the end of its spec. A word \emph{holds}
  $n \le 64$ bits of a pattern of period $m$ from a position on when its bit $t < n$ is the bit of
  the pattern at that position plus $t$ and its other bits are clear; the words of a row are the
  windows of the pattern when word $j$ holds its $64$ bits from $64 j$ on; and a row whose words are
  the windows of a pattern that a test reads at the residues of the numerators is the closed form of
  the test, which is how the contracts of the builders conclude. The row of a denominator in words
  (Definition~\ref{def:row}) of a sieve entry with a positive modulus is the closed form of the
  admission test of the entry.
\end{lemma}

\begin{proof}
  \leanok
  Two words are equal when their bits are; the bits of the closed form are those of the test, and
  the bit $64 j + t$ and the bit $64 (j \bmod m\,\mathrm{rbaPack}) + t$ differ by a multiple of
  $64\,\mathrm{rbaPack}\, m$, a multiple of $m$.
\end{proof}

A sieve spec gains the start and the end of the walk of the chunked arm through its row, indices
into the row for the pointers \verb|start| and \verb|end| of the \verb|sieve_spec| of the C code,
the end named $\mathrm{stop}$ since \verb|end| is a keyword of Lean (Definition~\ref{def:rows}); a
sieve spec \emph{fits} when, besides the bounds of its modulus and its offset, its end is $p$ and
its row wraps (Lemma~\ref{lem:rowIndices}), which the sieve specs of a denominator do, their rows
being closed forms (Section~\ref{sec:storeRun}).

\section{The chunked first stage}\label{sec:chunked}

\begin{definition}[Bit arrays loaded, stored and filled]\label{def:bitArrayMoves}
  \uses{def:rbaPack}
  \lean{Ratpoints.Sieve.loadBitArray, Ratpoints.Sieve.storeBitArray, Ratpoints.Sieve.fillBitArray, Ratpoints.Sieve.getElem!_loadBitArray, Ratpoints.Sieve.getElem!_storeBitArray, Ratpoints.Sieve.getElem!_fillBitArray}
  \leanok
  Next to the and of two bit arrays (Definition~\ref{def:bitArrayOps}), one word at a time: bit
  array $j$ of an array of words set to the class word, the 2-adic pattern of the class of the
  denominator in one word (each word of \verb|bits64|), and-ed with bit array $x$ of a row, the load
  \verb|reg = bits64 & *siv| of a register; bit array $y$ set to bit array $j$ of the registers, the
  store \verb|surv[y] = reg|; and bit array $y$ set to the class word, the store
  \verb|survivors[y] = bits64| of both arms. They are expressions inside the lines of their callers
  and have no C ranges of their own.
\end{definition}

\begin{definition}[The starts of the rows]\label{def:chunkStarts}
  \uses{def:rows, def:modMul, def:useCLong}
  \lean{Ratpoints.Sieve.chunkStarts, Ratpoints.Sieve.wrapStart}
  \leanok
  The sieve specs of the first stage with their starts at the first word number $w_0$ of a piece, in
  a local array: for each modulus $p$ with its offset and its reciprocal, stepped together, the
  start $(w_0 + \mathrm{offset}) \bmod p$, by \verb|mod_mul| when \verb|small| holds and by
  \verb|mod| otherwise (\csource{sift.c}{596}{606}); and a start stepped back by the modulus while
  it is at least the end, \verb|while (siv >= end) siv -= p| (\csource{sift.c}{683}{683},
  \csource{sift.c}{805}{805}).
\end{definition}

\begin{definition}[The whole chunks]\label{def:chunkSieve}
  \uses{def:chunkStarts, def:bitArrayMoves, def:bitArrayOps, def:chunk}
  \lean{Ratpoints.Sieve.chunkSieve, Ratpoints.Sieve.andChunks}
  \leanok
  One chunk of $\mathrm{chunk}$ bit arrays from the starts of the local sieve specs: the registers,
  a vector of $\mathrm{chunk}$ bit arrays, loaded with the class word and-ed with the bit arrays
  $\mathrm{start}, \dots, \mathrm{start} + \mathrm{chunk} - 1$ of the row of the first modulus,
  whose start moves on by $\mathrm{chunk}$ and is stepped back; then for each further modulus the
  registers and-ed with its bit arrays from its start on, and its start moved in the same way; the
  registers and the specs with the new starts returned (\csource{sift.c}{625}{681},
  \csource{sift.c}{684}{804}, \csource{sift.c}{806}{863}, with both blocks of output under
  \verb|DEBUG|, which the model does not print). The whole chunks of a piece, the bit arrays below
  $\mathrm{full} = n - n \bmod \mathrm{chunk}$: for
  $c = 0, \mathrm{chunk}, 2\,\mathrm{chunk}, \dots$ below $\mathrm{full}$, one chunk, its registers
  stored into the bit arrays $c, \dots, c + \mathrm{chunk} - 1$ of the piece
  (\csource{sift.c}{624}{624}, \csource{sift.c}{864}{913}).
\end{definition}

\begin{definition}[The tail]\label{def:tailLegs}
  \uses{def:bitArrayMoves, def:bitArrayOps, def:chunk, def:rows}
  \lean{Ratpoints.Sieve.tailLeg, Ratpoints.Sieve.tailLegs}
  \leanok
  One leg of $w \le 8$ registers at the offset $\mathrm{off}$ into the tail, \verb|tail_leg|: a
  local vector of $8$ bit arrays loaded with the class word and-ed with the bit arrays
  $\mathrm{start} + \mathrm{off}, \dots$ of the row of the first modulus, and-ed with those of the
  further moduli, and stored into the bit arrays $\mathrm{full} + \mathrm{off}, \dots$ of the piece,
  the starts left as they are (\csource{sift.c}{452}{465}, \csource{sift.c}{473}{496}). The tail of
  $t = n - \mathrm{full}$ bit arrays in the four calls of the C code, each under its guard:
  $8 < \mathrm{chunk}$ and $t \mathbin{\&} 8 \ne 0$ for the leg of $8$ at $0$, $4 < \mathrm{chunk}$
  and $t \mathbin{\&} 4 \ne 0$ for the leg of $4$ at $t \mathbin{\&} 8$, $2 < \mathrm{chunk}$ and
  $t \mathbin{\&} 2 \ne 0$ for the leg of $2$ at $t \mathbin{\&} 12$, and $t \mathbin{\&} 1 \ne 0$
  for the leg of $1$ at $t \mathbin{\&} 14$ (\csource{sift.c}{915}{940}, with the dead
  \verb|#error|).
\end{definition}

\begin{definition}[The chunked first stage]\label{def:andRowsChunked}
  \uses{def:chunkStarts, def:chunkSieve, def:tailLegs, def:bitArrayMoves}
  \lean{Ratpoints.Sieve.andRowsChunked}
  \leanok
  The first stage over a piece of $n$ bit arrays from the word number $w_0$, the arm of
  \verb|_ratpoints_sift0| for a chunk from $2$ to $16$, with the arguments of the first stage one
  row at a time (Definition~\ref{def:andRows}): the local sieve specs with their starts; with no
  modulus, every bit array set to the class word, from the last; otherwise the whole chunks, then
  the tail from the starts after them (\csource{sift.c}{575}{575}, \csource{sift.c}{608}{620},
  \csource{sift.c}{941}{941}; with the constructor \verb|RBA| of the class word, which both arms
  share, \csource{find_points.c}{1098}{1098}, \csource{rp-private.h}{129}{132},
  \csource{rp-private.h}{203}{204}, \csource{rp-private.h}{235}{235},
  \csource{rp-private.h}{261}{261}, \csource{rp-private.h}{285}{285}). The rows must have their
  copies, which every row of the model has (Section~\ref{sec:storeRun}).
\end{definition}

\begin{lemma}[The chunked first stage and-s the rows]\label{lem:andRowsChunked}
  \uses{def:andRowsChunked, def:closedRow, lem:rowIndices, def:andRows}
  \lean{Ratpoints.Sieve.SieveSpec.startAt, Ratpoints.Sieve.SieveSpec.advance, Ratpoints.Sieve.SieveSpec.WalkFits, Ratpoints.Sieve.SieveSpec.Fits.walkFits_startAt, Ratpoints.Sieve.partialReg, Ratpoints.Sieve.andFrom, Ratpoints.Sieve.ChunkStartsInvariant, Ratpoints.Sieve.chunkStarts_spec, Ratpoints.Sieve.WrapStartInvariant, Ratpoints.Sieve.wrapStart_spec, Ratpoints.Sieve.ChunkModInvariant, Ratpoints.Sieve.ChunkAndInvariant, Ratpoints.Sieve.chunkSieve_spec, Ratpoints.Sieve.AndChunksInvariant, Ratpoints.Sieve.andChunks_spec, Ratpoints.Sieve.LegAndInvariant, Ratpoints.Sieve.LegStoreInvariant, Ratpoints.Sieve.tailLeg_spec, Ratpoints.Sieve.tail_bits, Ratpoints.Sieve.TailDone, Ratpoints.Sieve.tailLegs_spec, Ratpoints.Sieve.andRowsChunked_spec, Ratpoints.Sieve.andRowsChunked_post}
  \leanok
  For sieve specs that fit (Lemma~\ref{lem:rowIndices}), the reciprocals of their moduli and the
  bounds of \verb|small|: the local array holds the specs with their starts at $w_0$, the remainder
  of $w_0 + \mathrm{offset}$ (Lemma~\ref{lem:rowIndices}); for the end $p > 0$ a start stepped back
  is its remainder modulo $p$; a chunk loads into its registers the class word and-ed with the rows
  of all the moduli from their starts on and moves every start on by the chunk modulo $p$; the whole
  chunks below a multiple $\mathrm{full}$ of the chunk store into their bit arrays the words that
  the stage reads from the starts and move the starts on by $\mathrm{full}$, for specs with a
  positive modulus, the end $p$, a start below $p$ and a row that wraps; a leg of $w \le 8$
  registers stores into its bit arrays the class word and-ed with the rows at the starts plus its
  place in the tail; and the legs of a tail of $t < \mathrm{chunk} \le 16$ bit arrays do its bit
  arrays $\mathrm{full}, \dots, n - 1$. So, for $\mathrm{chunk} \le 16$, the chunked first stage has
  the postcondition of the first stage one row at a time (Definition~\ref{def:andRows}): word
  $y < n\,\mathrm{rbaPack}$ of the piece is the class word and-ed with the row words that the specs
  read for it at their offsets.
\end{lemma}

\begin{proof}
  \leanok
  The start of the chunk at $c$ is $(w_0 + \mathrm{offset} + c) \bmod p < p$, and the chunk reads
  the bit arrays from it up to $\mathrm{chunk} - 1$ further on, below $p + \mathrm{chunk} - 1$,
  where the row wraps: bit array $s + j$ is bit array $(s + j) \bmod p = (w_0 + \mathrm{offset} + c
  + j) \bmod p$, the one that the first stage one row at a time reads for bit array $c + j$. The
  bits of $t < 16$ are decided over the sixteen values: each leg that runs has its width, at most
  $t$, and its offset is the sum of the widths of the wider legs, so that the legs cover $0, \dots,
  t - 1$ once each; the guards on the chunk drop only legs whose bit is clear for $t <
  \mathrm{chunk}$. A chunk above $16$ would leave a tail of $16$ or more bit arrays, which the legs
  do not cover, so the contract asks for $\mathrm{chunk} \le 16$, which the test of the arm gives.
  The loops end by measures, the step back by the start and the chunks by $\mathrm{full} - c$.
\end{proof}

The first stage one row at a time (Definition~\ref{def:andRows}) writes the class word into the bit
arrays from the last, as the C code does, and one piece (Definition~\ref{def:siftPiece}) chooses the
arm by the test $1 < \mathrm{chunk} \le 16$ of the conditional compilation of the C code
(\csource{sift.c}{452}{452}, \csource{sift.c}{575}{575}, \csource{sift.c}{943}{943}): the chunked
arm for the chunk $16$ of the model, and the other arm for a chunk of $1$ or above $16$. Its
contract keeps its statement, the first stage of either arm giving the same words.

\section{The tables of a prime power as words}\label{sec:powerWords}

\begin{definition}[Bit arrays of words for the tables of a prime power]\label{def:powerWords}
  \uses{def:barrett, def:modinv}
  Proposed names: \lname{Ratpoints.modWords}, \lname{Ratpoints.Sieve.bitAt},
  \lname{Ratpoints.Sieve.setBitAt}, \lname{Ratpoints.Sieve.patBit},
  \lname{Ratpoints.Sieve.setSquareBits}, \lname{Ratpoints.Sieve.unitInverses}. The number of words
  of a pattern of $\mathrm{maxPrimeEven}$ bits, $\mathrm{modWords} = \lceil \mathrm{maxPrimeEven}/64
  \rceil$ (\verb|RP_MODWORDS|, \csource{rp-private.h}{523}{523}); bit $x$ of a bit array of words,
  bit $x \bmod 64$ of word $\lfloor x/64 \rfloor$, read and set as the C code reads and sets the
  bits of the tables of a prime power and of the patterns of the rows, and the same bit read off the
  bits of the word, as the proofs read the patterns; the bits of the squares $x^2 \bmod m$ of the
  residues $x < m$ set in the bit array \verb|sq| (\csource{find_points.c}{2283}{2284}); and the
  inverses of the units modulo $m = p^e$, $0$ for the other residues, in the loop over the residues
  (\csource{find_points.c}{2302}{2303}). The tables of a prime power
  (Definition~\ref{def:examinePower}) keep \verb|fsq| and \verb|gsq| as bit arrays of
  $\mathrm{modWords}$ words, which \verb|examine_power| clears together with \verb|sq| and sets bit
  by bit, and its admission test reads their bits.
\end{definition}

\begin{lemma}[The bit arrays of a prime power]\label{lem:powerWords}
  \uses{def:powerWords, lem:modinv}
  Proposed names: \lname{Ratpoints.Sieve.maxPrimeEven\_le\_sixty\_four\_mul\_modWords},
  \lname{Ratpoints.Sieve.SquareBitsInvariant}, \lname{Ratpoints.Sieve.setSquareBits\_spec},
  \lname{Ratpoints.Sieve.UnitInversesInvariant}, \lname{Ratpoints.Sieve.unitInverses\_spec},
  \lname{Ratpoints.Sieve.FsqInvariant}, \lname{Ratpoints.Sieve.GsqInvariant}. A modulus below
  $\mathrm{maxPrimeEven}$ fits the $\mathrm{modWords}$ words, $\mathrm{maxPrimeEven} \le
  64\,\mathrm{modWords}$; on a bit array of $\mathrm{modWords}$ words and $m \le
  64\,\mathrm{modWords}$ the bits of the squares modulo $m$ get set; and for a prime $p$ and $e \ge
  1$ the inverses are below $p^e$, that of a unit its inverse modulo $p^e$. With the invariants of
  the loops of \verb|fsq| and \verb|gsq|, the contract of the tables of a prime power and the lemma
  that its rows let every point through are restated on the bits (Lemma~\ref{lem:power_admits}), for
  $p^e \le 64\, \mathrm{modWords}$.
\end{lemma}

\begin{proof}
  The bit set by \verb|setBitAt| is the bit read by \verb|bitAt| and no other bit changes; the loop
  over the multiples of $p$ is a \verb|while| loop with the measure $m - t$; the inverses are those
  of the extended Euclidean algorithm (Lemma~\ref{lem:modinv}).
\end{proof}

\section{The rows of a prime}\label{sec:primeRows}

\begin{definition}[The words of a forbidden divisor and the row of $0$]\label{def:sieves0}
  \uses{def:chunk, def:rbaPack}
  Proposed names: \lname{Ratpoints.Sieve.forbiddenRow}, \lname{Ratpoints.Sieve.sieves0}. The $p$
  words of the forbidden divisor $p$, with the bits of the multiples $p i$, $i < 64$, cleared: the
  array \verb|work| of the generator of \verb|find_points.h| for the primes of the table, and the
  words that \verb|sieving_info| makes for the forbidden divisors beyond it
  (\csource{find_points.c}{2773}{2777}, \csource{gen_find_points_h.c}{118}{120}); they move here
  from the forbidden divisors (Definition~\ref{def:forbidden}), which read them. And for a prime $p$
  of the table these words written $\mathrm{rbaPack}$ times one after the other, then the copies,
  word $i$ of $\mathrm{rbaPack}\,(\mathrm{chunk} - 1)$ of them the word $i \bmod p$: the row of $p$
  in the table \verb|sieves0| that the build generates and \verb|find_points.c| includes
  (\csource{find_points.c}{37}{37}, \csource{gen_find_points_h.c}{102}{116},
  \csource{gen_find_points_h.c}{122}{134}), a function of $p$ as the squares and the inverses
  modulo $p$ are.
\end{definition}

\begin{lemma}[The row of $0$]\label{lem:sieves0}
  \uses{def:sieves0, def:closedRow, lem:closedRow, def:admits}
  Proposed names: \lname{Ratpoints.Sieve.getLsbD\_getElem!\_forbiddenRow},
  \lname{Ratpoints.Sieve.size\_forbiddenRow}, \lname{Ratpoints.Sieve.size\_sieves0},
  \lname{Ratpoints.Sieve.getElem!\_sieves0}, \lname{Ratpoints.Sieve.sieves0\_eq\_closedRow},
  \lname{Ratpoints.Sieve.extract\_sieves0}. Bit $t$ of word $r < p$ of the words of the forbidden
  divisor $p$ is clear exactly when $p \mid 64 r + t$, the lemma that moves here from the forbidden
  divisors (Definition~\ref{def:forbidden}) into a file without Mathlib; for $p > 0$ the row of $0$
  is the closed form of the row of $0$ of the table of $p$, whatever its entries: bit $t$ of word
  $x$ set exactly when $p \nmid 64 x + t$; and its first $p$ words are the words of the forbidden
  divisor $p$.
\end{lemma}

\begin{proof}
  The $64$ multiples $p i$, $i < 64$, are the multiples of $p$ below $64 p$, the bits of the $p$
  words; word $j$ of the row is word $j \bmod p$ of the words of the forbidden divisor, and
  $64 j + t \equiv 64 (j \bmod p) + t \pmod p$.
\end{proof}

\begin{definition}[The kinds of builders]\label{def:rowInit}
  Proposed names: \lname{Ratpoints.Sieve.RowInit}, \lname{Ratpoints.Sieve.sieveInit}. The builder of
  the rows of a sieve entry by its kind, for the pointer \verb|init| of the type
  \verb|ratpoints_init_fun| (\csource{rp-private.h}{497}{498}): the function that the build
  generates for a prime below $64$ from \verb|CODE_INIT_SIEVE1|, that for a prime above $64$ from
  \verb|CODE_INIT_SIEVE2|, the builder of the rows of a prime power and that of the rows of a
  product. The kind of the builder of a prime $p$ of the table: that from \verb|CODE_INIT_SIEVE1|
  for $p < 64$ and that from \verb|CODE_INIT_SIEVE2| otherwise, the entry of $p$ in the table
  \verb|sieve_init| of pointers that the build generates (\csource{find_points.c}{133}{133},
  \csource{gen_init_sieve_h.c}{48}{62}, \csource{init.c}{246}{246}). The sieve entries hold the kind
  of their builder, with the reference to their table of rows (Definition~\ref{def:SieveEntry}).
\end{definition}

\begin{definition}[The macros of the builders of a prime]\label{def:initMacros}
  \uses{def:chunk, def:rbaPack}
  Proposed names: \lname{Ratpoints.Sieve.initBit}, \lname{Ratpoints.Sieve.initFourDecl},
  \lname{Ratpoints.Sieve.initFourLoop}, \lname{Ratpoints.Sieve.initReplicate}. One definition for
  each macro whose body is a unit of code, in the form of the production builds: a bit of a word set
  from an entry of the table by a shift, \verb|RP_INIT_BIT| (\csource{init.c}{42}{46}); the four
  residues $d$, $2d$, $3d$ and the step $4d$ modulo $p$, \verb|RP_INIT_FOUR_DECL|
  (\csource{init.c}{62}{63}, \csource{init.c}{65}{70}); the loop that sets four bits of a word at a
  time from the four residues, stepping each by $4d \bmod p$, \verb|RP_INIT_FOUR_LOOP|
  (\csource{init.c}{64}{64}, \csource{init.c}{72}{83}); and the copies of the first $p$ words of a
  row into the next $(\mathrm{rbaPack} - 1)\, p$, so that word $j$ is word $j \bmod p$, and of the
  first $(\mathrm{chunk} - 1)\,\mathrm{rbaPack}$ words to the end, \verb|RP_INIT_REPLICATE|
  (\csource{init.c}{99}{111}), whose forward copies read words already copied, so that a prime below
  $\mathrm{chunk} - 1$ is copied as in the C code.
\end{definition}

\begin{definition}[The rows of a prime below $64$]\label{def:sieveInit1}
  \uses{def:initMacros, def:entryRef, def:examinePrime}
  Proposed names: \lname{Ratpoints.Sieve.sieve1Pattern}, \lname{Ratpoints.Sieve.sieve1Rows},
  \lname{Ratpoints.Sieve.sieveInit1}. For the table of a prime $p < 64$
  (Definition~\ref{def:examinePrime}), the reference to its entry, a residue $b$ and the store of
  the tables of rows (Definition~\ref{def:entryRef}), as the function that the build generates for
  $p$ from \verb|CODE_INIT_SIEVE1|: the pattern of $p$ bits in a word, bit $a$ the entry of the
  table of the prime at $a\, d \bmod p$ for $d$ the inverse of $b$ from it, by the loop of four bits
  and then the plain loop (\csource{init.c}{131}{143}); the pattern shifted right by $64 \bmod p$
  for the second word, and the pattern repeated through the word by or-ing in itself shifted left by
  $i = p, 2p, 4p, \dots$ below $64$; the first $p$ words of the row rotated out of the two words, by
  the shifts $\mathrm{diff\_shift} = (\lfloor 64/p \rfloor + 1)\, p \bmod 64$ and $64 -
  \mathrm{diff\_shift}$ (\csource{init.c}{160}{177}); the copies; and the row stored into the table
  of rows of the entry, \verb|RP_PUBLISH| (\csource{init.c}{118}{130}, \csource{init.c}{144}{159},
  \csource{init.c}{178}{183}).
\end{definition}

\begin{definition}[The rows of a prime above $64$]\label{def:sieveInit2}
  \uses{def:initMacros, def:entryRef, def:examinePrime}
  Proposed names: \lname{Ratpoints.Sieve.sieve2Word}, \lname{Ratpoints.Sieve.sieve2Pattern},
  \lname{Ratpoints.Sieve.sieve2Rows}, \lname{Ratpoints.Sieve.sieveInit2}. For the table of a prime
  $p > 64$ (Definition~\ref{def:examinePrime}), the reference to its entry, a residue $b$ and the
  store of the tables of rows, as the function that the build generates for $p$ from
  \verb|CODE_INIT_SIEVE2|: the $\lfloor p/64 \rfloor + 2$ words \verb|help| cleared from the last;
  word $n \le \lfloor p/64 \rfloor$ of the pattern, $64$ bits for $n < \lfloor p/64 \rfloor$ and $p
  \bmod 64$ for the last, by the loop of four bits and the plain loop, the four residues carried
  from word to word (\csource{init.c}{209}{216}, \csource{init.c}{204}{208},
  \csource{init.c}{217}{219}); the first $p$ words of the row, the first $\lfloor p/64 \rfloor$
  copied and the others rotated out of the ring of $\lfloor p/64 \rfloor + 1$ words by the shifts $p
  \bmod 64$ and $64 - (p \bmod 64)$ (\csource{init.c}{220}{238}); the copies; and the row stored
  into the table of rows of the entry (\csource{init.c}{186}{203}, \csource{init.c}{239}{244}).
\end{definition}

\begin{lemma}[The builders of a prime build the closed forms]\label{lem:primeRows}
  \uses{def:sieveInit1, def:sieveInit2, lem:closedRow, def:admits}
  Proposed names: \lname{Ratpoints.Sieve.lanePattern}, \lname{Ratpoints.Sieve.LanesAt},
  \lname{Ratpoints.Sieve.FourLanes}, \lname{Ratpoints.Sieve.initFourDecl\_lanes},
  \lname{Ratpoints.Sieve.FourInvariant}, \lname{Ratpoints.Sieve.initFourLoop\_spec},
  \lname{Ratpoints.Sieve.PatInvariant}, \lname{Ratpoints.Sieve.sieve1Pattern\_spec},
  \lname{Ratpoints.Sieve.PatternDoubleInvariant}, \lname{Ratpoints.Sieve.RotInvariant},
  \lname{Ratpoints.Sieve.sieve1Rows\_spec}, \lname{Ratpoints.Sieve.sieve1\_shifts},
  \lname{Ratpoints.Sieve.RepInvariant}, \lname{Ratpoints.Sieve.CopyInvariant},
  \lname{Ratpoints.Sieve.initReplicate\_spec}, \lname{Ratpoints.Sieve.sieve2Word\_spec},
  \lname{Ratpoints.Sieve.sieve2Pattern\_spec}, \lname{Ratpoints.Sieve.PatternWords},
  \lname{Ratpoints.Sieve.RingState}, \lname{Ratpoints.Sieve.sieve2Rows\_spec},
  \lname{Ratpoints.Sieve.sieve2\_shifts},
  \lname{Ratpoints.Sieve.toArray\_eq\_closedRow\_of\_rowWords},
  \lname{Ratpoints.Sieve.sieveInit1\_spec}, \lname{Ratpoints.Sieve.sieveInit2\_spec}. For the table
  of a prime $p$ and a residue $0 < b < p$ whose inverse is below $p$: for $p < 64$ the builder from
  \verb|CODE_INIT_SIEVE1| and for $64 \le p$ with $p \bmod 64 \ne 0$ that from
  \verb|CODE_INIT_SIEVE2| return the closed form of the row of $b$ of the table of the prime
  (Definition~\ref{def:closedRow}) and the store with that row in the table of rows of the entry.
  The pattern of the four residues is the pattern of the table of the prime at $a\, d \bmod p$, the
  copies make word $j$ of the row word $j \bmod p$, and every shift that the builders execute lies
  in $[0, 63]$, where those of the C code are defined: for $p < 64$, $\mathrm{diff\_shift}$ lies in
  $[1, 63]$ for every $p$, and for $p > 64$ the shift $p \bmod 64$ is not $0$.
\end{lemma}

\begin{proof}
  The four residues are $(a + i)\, d \bmod p$ for $i < 4$ at the bit $a$, stepped by $4d \bmod p$;
  the invariant of each loop says which bits of the pattern the word holds. The doubling repeats the
  pattern with the period $p$ through $64$ bits; the rotation keeps in the two words, or in the
  ring, the windows of the periodic pattern at $64 a$ and after it; the windows of the pattern of
  the table of the prime read at $N\, d \bmod p$ are the closed form of the admission test at $b \ne
  0$, which reads that table at $(N \bmod p)\, d \bmod p$. At $b = 0$ the builders would give the
  wrong row (every lane at the residue $0$), which is why the row of $0$ is \verb|sieves0| and the
  builders are called at $b \ne 0$ only (Lemma~\ref{lem:fetchRow}).
\end{proof}

\section{The rows of a prime power and of a product}\label{sec:compositeRows}

\begin{definition}[The layout of a pattern]\label{def:layOutPattern}
  \uses{def:powerWords, def:chunk, def:rbaPack}
  Proposed names: \lname{Ratpoints.Sieve.window}, \lname{Ratpoints.Sieve.repeatPattern},
  \lname{Ratpoints.Sieve.layOutPattern}. The $64$ bits of an array of words from a position on, cut
  out of two words by two shifts, for line 291 of \verb|lay_out_pattern|; the pattern of period $m$
  of an array of words repeated through $\mathrm{words} = \lfloor (m + 127)/64 \rfloor$ words, with
  $64\,\mathrm{words} \ge m + 64$, the words cleared first and each copy or-ed in at its offset
  $o = 0, m, 2m, \dots$, word by word and with the guard $s \ne 0$ of the C code before the shift by
  $64 - s$ (\csource{init.c}{268}{279}); and the row of $m + \mathrm{chunk} - 1$ bit arrays: for
  $x < m\,\mathrm{rbaPack}$, word $x$ the $64$ bits of the repeated pattern from the offset
  $o = 64 x \bmod m$ on, the offset stepped by $64 \bmod m$ modulo $m$, then the copies of the first
  $(\mathrm{chunk} - 1)\,\mathrm{rbaPack}$ words, with its own loop of copies
  (\csource{init.c}{259}{267}, \csource{init.c}{280}{297}).
\end{definition}

\begin{definition}[The rows of a prime power]\label{def:sieveInitPower}
  \uses{def:layOutPattern, def:powerWords, def:examinePower, def:entryRef}
  Proposed names: \lname{Ratpoints.Sieve.powerPatternUnit},
  \lname{Ratpoints.Sieve.powerPatternNonunit}, \lname{Ratpoints.Sieve.sieveInitPower}. For the
  tables of a prime power $m = p^e$ (Definition~\ref{def:examinePower}), the reference to its entry,
  a residue $b$ and the store of the tables of rows, as \verb|_ratpoints_sieve_init_power| builds
  the row: the pattern of $\mathrm{modWords}$ words cleared; for a unit $b$, bit $x$ set when bit
  $r$ of \verb|fsq| is, the residue $r$ stepped by the inverse of $b$ modulo $m$
  (\csource{init.c}{369}{375}); for $p \mid b$, bit $x$ of a unit $x$ set when bit $b\, x^{-1} \bmod
  m$ of \verb|gsq| is, the other residues skipped by \verb|continue| (\csource{init.c}{378}{385});
  the pattern laid out as a row (Definition~\ref{def:layOutPattern}); and the row stored into the
  table of rows of the entry (\csource{init.c}{359}{368}, \csource{init.c}{376}{377},
  \csource{init.c}{386}{387}, \csource{init.c}{391}{393}, with its prototype
  \csource{rp-private.h}{592}{592}).
\end{definition}

\begin{definition}[The rows of a product]\label{def:sieveInitProduct}
  \uses{def:entryRef, def:rowInit, def:sieveInit1, def:sieveInit2, def:sieveInitPower,
    def:SieveEntry, def:bitArrayMoves, def:bitArrayOps}
  Proposed names: \lname{Ratpoints.Sieve.initFactorRow}, \lname{Ratpoints.Sieve.factorRows},
  \lname{Ratpoints.Sieve.productCopy}, \lname{Ratpoints.Sieve.productAnd},
  \lname{Ratpoints.Sieve.sieveInitProduct}. The call \verb|(*(fe->init))(fe, bf, args)| through the
  pointer of the entry of a prime or a prime power: a match on the kind of builder
  (Definition~\ref{def:rowInit}) that calls the builder of that kind on the table of the prime or
  the tables of the prime power of the factor, with the reference given, which stores the row into
  that table of rows (a kind that does not fit the factor, which no entry of the model has, gives
  the empty row). For the sieve entry of a product $m$, a residue $b$ and the store, as
  \verb|_ratpoints_sieve_init_product| builds the row: for each factor of modulus $q$, the row of $b
  \bmod q$ from the table of rows of the entry of the factor, and when it is absent the row that the
  builder of that entry builds into it (\csource{init.c}{407}{414}); the row of the first factor
  laid down in blocks of its $q$ bit arrays (\csource{init.c}{419}{423}); the rows of the others
  and-ed in, in blocks of theirs (\csource{init.c}{424}{429}); the copies of the first
  $\mathrm{chunk} - 1$ bit arrays after the first $m$; and the row stored into the table of rows of
  the entry of the product (\csource{init.c}{399}{406}, \csource{init.c}{415}{418},
  \csource{init.c}{430}{430}, \csource{init.c}{443}{445}, with its prototype
  \csource{rp-private.h}{593}{593}).
\end{definition}

\begin{lemma}[The builders of a prime power and of a product]\label{lem:compositeRows}
  \uses{def:layOutPattern, def:sieveInitPower, def:sieveInitProduct, lem:closedRow}
  Proposed names: \lname{Ratpoints.Sieve.getLsbD\_window}, \lname{Ratpoints.Sieve.PatClear},
  \lname{Ratpoints.Sieve.PatZero}, \lname{Ratpoints.Sieve.repeatPattern\_spec},
  \lname{Ratpoints.Sieve.LayOuterInvariant}, \lname{Ratpoints.Sieve.LayInnerInvariant},
  \lname{Ratpoints.Sieve.RowLoopInvariant}, \lname{Ratpoints.Sieve.layOutPattern\_spec},
  \lname{Ratpoints.Sieve.powerPatternUnit\_spec}, \lname{Ratpoints.Sieve.powerPatternNonunit\_spec},
  \lname{Ratpoints.Sieve.sieveInitPower\_spec}, \lname{Ratpoints.Sieve.factorWord},
  \lname{Ratpoints.Sieve.productWord}, \lname{Ratpoints.Sieve.FirstInvariant},
  \lname{Ratpoints.Sieve.FactorInvariant}, \lname{Ratpoints.Sieve.productCopy\_spec},
  \lname{Ratpoints.Sieve.productAnd\_spec},
  \lname{Ratpoints.Sieve.toArray\_eq\_closedRow\_of\_productWord},
  \lname{Ratpoints.Sieve.Factor.admits\_mod\_of\_dvd},
  \lname{Ratpoints.Sieve.SieveEntry.admits\_eq\_all}. For a pattern of period $0 < m \le
  64\,\mathrm{modWords}$ without bits from $m$ on, the layout gives the row whose words are the
  windows of the pattern; for the tables of a prime power $m \le 64\,\mathrm{modWords}$ and a
  residue $b < m$ whose inverse in the table is below $m$, the pattern for a unit $b$ has bit $x$
  the bit of \verb|fsq| at $x\, b^{-1} \bmod m$, that for $p \mid b$ bit $x$ the bit of \verb|gsq|
  at $b\, x^{-1} \bmod m$ for a unit $x$, and the builder returns the closed form of the row of $b$
  of the tables of the prime power, and the store with it in the table of rows of the entry; for a
  product with at least one factor, whose factors have positive moduli dividing $m$, the first
  factor laid down and the others and-ed in give word $y < m\,\mathrm{rbaPack}$ the and of the words
  of the factors read in their blocks, and from rows of the factors that are the closed forms of
  their rows at $b \bmod q$ the row, repeated after its first $m$ bit arrays, is the closed form of
  the row of $b$ of the product.
\end{lemma}

\begin{proof}
  \uses{lem:admits_mul}
  The repeated pattern holds the pattern at every bit up to $64\,\mathrm{words}$, so the window at
  $64 x \bmod m$ is the window at $64 x$ of the periodic pattern; the guard $s \ne 0$ keeps every
  shift in $[1, 63]$. The admission test of a prime power reads \verb|fsq| at $x\, r^{-1}$ or
  \verb|gsq| at $r\, x^{-1}$, which are the bits that the patterns set. The moduli of the factors
  divide $m$, so a block of $q$ bit arrays at a multiple of $q$ is the row of the factor, and the
  admission test of a factor reads the numerators modulo its modulus.
\end{proof}

\section{The tables of rows}\label{sec:rowStore}

\begin{definition}[The tables of rows]\label{def:entryRef}
  Proposed names: \lname{Ratpoints.Sieve.EntryRef}, \lname{Ratpoints.Sieve.RowStore},
  \lname{Ratpoints.Sieve.RowStore.row?}, \lname{Ratpoints.Sieve.RowStore.publish}. A reference to
  the sieve entry of ratpoints that holds a table of rows, for a pointer to the entry: a prime $p$
  by itself, for the entry \verb|prime_se[pn]| of $p = \mathrm{prime}[\mathrm{pn}]$ that
  \verb|examine_prime| makes; a prime power $p^e$ by its prime and its exponent, for the entry
  \verb|power_se[idx]| that \verb|make_power| makes once and that the prime power as a modulus and
  as a factor of products share, whose index \verb|add_moduli| keeps in \verb|pw_idx[pn][e]|; and a
  product by the position $n$ of its modulus among the moduli of the parameters, for the entry of
  its own that \verb|make_modulus| makes. The store of the run maps each reference to its table of
  rows, the rows by the residue, $m$ of them for the modulus $m$, each absent until it is built, for
  the arrays \verb|sieve| of the entries with \verb|NULL| for an absent row; the row of a residue in
  it, none for an absent row or table of rows; and the store with a row put into a table of rows,
  which each builder does at its end, as \verb|RP_PUBLISH| stores the row into the entry, and which
  leaves a store without that table of rows as it is.
\end{definition}

\begin{definition}[The row of a sieve entry, read or built]\label{def:fetchRow}
  \uses{def:sieveInitProduct, def:rowInit, def:entryRef, def:SieveEntry}
  Proposed names: \lname{Ratpoints.Sieve.SieveEntry.initRow},
  \lname{Ratpoints.Sieve.SieveEntry.fetchRow}. The call \verb|(*(se->init))(se, b, args)| of the run
  through the pointer of a sieve entry: for a product the builder of the rows of a product, which
  builds the absent rows of its factors into their tables of rows first, and for a prime or a prime
  power the builder of its kind on the table of the prime or the tables of the prime power of its
  one factor, with the reference to the entry (Definition~\ref{def:sieveInitProduct}). And the row
  of a residue $b < p$ of a sieve entry: the row in the table of rows of its entry when it is there,
  and otherwise the row that its builder builds into that table of rows, the body of the loop of
  \verb|sift| that builds the missing rows, which reads the pointer to the row again before it calls
  the builder (Definition~\ref{def:buildRows}).
\end{definition}

\begin{lemma}[Rows read or built are closed forms]\label{lem:fetchRow}
  \uses{def:fetchRow, lem:primeRows, lem:compositeRows, def:closedRow}
  Proposed names: \lname{Ratpoints.Sieve.RowStore.Grows},
  \lname{Ratpoints.Sieve.RowStore.row?\_publish}, \lname{Ratpoints.Sieve.Factor.Buildable},
  \lname{Ratpoints.Sieve.Factor.Builds}, \lname{Ratpoints.Sieve.initFactorRow\_spec},
  \lname{Ratpoints.Sieve.FactorRow}, \lname{Ratpoints.Sieve.FactorRowsPre},
  \lname{Ratpoints.Sieve.FactorRowsInvariant}, \lname{Ratpoints.Sieve.factorRows\_spec},
  \lname{Ratpoints.Sieve.ProductRowsPre}, \lname{Ratpoints.Sieve.sieveInitProduct\_spec},
  \lname{Ratpoints.Sieve.InitRowPre}, \lname{Ratpoints.Sieve.SieveEntry.initRow\_spec},
  \lname{Ratpoints.Sieve.SieveEntry.fetchRow\_spec}. A factor is \emph{buildable} when its builder
  has what it needs: a prime $p > 0$ with the inverses of its residues below $p$ and $p \bmod 64 \ne
  0$ when $p \ge 64$, or a prime power $0 < m \le 64\,\mathrm{modWords}$ with its inverses below
  $m$; its builder is called at a residue below the modulus, other than $0$ for a prime. Then the
  call through the pointer gives the closed form of the row of the factor and puts it into the table
  of rows given. For a product with a positive modulus $m$ and at least one factor, whose factors
  are buildable, have positive moduli dividing $m$, rows in their tables of rows that are closed
  forms, an absent row only where their builders may be called, and one closed form for two factors
  with one reference, the rows of the factors that the builder of the product reads or builds are
  their closed forms, the store \emph{grows} by rows of the factors (the tables of rows and their
  sizes kept, no row replaced, the new rows those of the factors), and the row of the product is its
  closed form, put into its table of rows. The row of a sieve entry read or built is the closed form
  of its row, and the store is the same or has grown in this way, with the row put into the table of
  rows of the entry.
\end{lemma}

\begin{proof}
  \uses{lem:admits_mul}
  Lemmas~\ref{lem:primeRows} and~\ref{lem:compositeRows}, and for the product an invariant of the
  loop over its factors: each row read is a closed form by the hypothesis on the tables of rows,
  each row built is one by the contract of its builder, at $b \bmod q$, which for a prime is not $0$
  when its row is absent, its row of $0$ being in its table of rows from the start.
\end{proof}

\section{The tables of rows of a run}\label{sec:storeRun}

\begin{definition}[The tables of rows made with the entries]\label{def:storeTables}
  \uses{def:entryRef, def:sieves0, def:rowInit, def:SieveEntry}
  Proposed names: \lname{Ratpoints.Sieve.primeRows}, \lname{Ratpoints.Sieve.Factor.newRows},
  \lname{Ratpoints.Sieve.SieveEntry.addTables}. The table of rows of a prime: its row of $0$ the row
  of $0$ of Definition~\ref{def:sieves0}, the other $p - 1$ rows absent, as \verb|examine_prime|
  sets the array \verb|sieve| of the entry (\csource{find_points.c}{1820}{1821}); the table of rows
  that the entry of a factor starts with, that of its prime, or for a prime power $m$ the $m$ rows
  absent; and the store with the tables of rows of a sieve entry, each made only when its reference
  is absent: the table of rows of the entry of each factor, and for a product the $m$ rows absent of
  its own entry, as \verb|make_power| and \verb|make_modulus| clear the arrays \verb|sieve|
  (\csource{find_points.c}{2427}{2427}, \csource{find_points.c}{2471}{2471}). The loop of the tables
  (Definition~\ref{def:examinePrimes}) makes the table of rows of each examined prime that says
  something, as \verb|examine_prime| makes its entry, and returns the store; \verb|sieving_info|
  (Definition~\ref{def:sievingInfo}) makes each sieve entry at the position $n$, the number of the
  entries made before it (Definition~\ref{def:makeModulus}), and right after it the tables of rows
  of that entry, and returns the store with the entries.
\end{definition}

\begin{lemma}[The store fits the sieve entries]\label{lem:storeTables}
  \uses{def:storeTables, lem:fetchRow, lem:sieves0, lem:sieveEntry_admits, def:makeModulus,
    def:closedRow}
  Proposed names: \lname{Ratpoints.Sieve.RowStore.Holds}, \lname{Ratpoints.Sieve.StoreHolds},
  \lname{Ratpoints.Sieve.SieveEntry.RefFits},
  \lname{Ratpoints.Sieve.SieveEntry.refFits\_of\_refInit\_eq},
  \lname{Ratpoints.Sieve.SieveEntry.RowsFit}, \lname{Ratpoints.Sieve.RowStore.PrimeZeroRows},
  \lname{Ratpoints.Sieve.runRows}, \lname{Ratpoints.Sieve.RowStore.MadeFor},
  \lname{Ratpoints.Sieve.RowStore.madeFor\_nil}, \lname{Ratpoints.Sieve.RowStore.MadeFor.addTables},
  \lname{Ratpoints.Sieve.RowStore.MadeFor.storeHolds}, \lname{Ratpoints.Sieve.RowStore.Fits},
  \lname{Ratpoints.Sieve.PrimeStore}, \lname{Ratpoints.IsModulusEntry.rowsFit\_runRows}. A store
  \emph{holds} a function $R$ of the references and the residues when every row present at the
  residue $r$ of the table of rows of a reference is $R$ at them; it holds $R$ for a list of sieve
  entries when, besides, the table of rows of every prime factor of the entries has its row of $0$.
  A store is a \emph{store of the primes} when it holds every $R$ whose rows of $0$ of the primes
  are \verb|sieves0| and the table of rows of each prime in it has its row of $0$; the store of the
  loop of the tables is one (Lemma~\ref{lem:examinePrimes}). A sieve entry at the position $n$
  \emph{fits} $R$ when it has the reference and the kind of builder of its position and its factors,
  a positive modulus, at least one factor, moduli of the factors that divide its modulus, buildable
  factors, and $R$ at its reference and at the references of its factors the closed forms of their
  rows. A store \emph{fits} a list of sieve entries when it holds some $R$ for them that every entry
  fits at its position. For a run of the polynomial with the coefficients $c$ of degree $d$ and the
  sieve entries $\mathrm{es}$, the rows of the references are: for a prime $p$ the closed forms of
  the rows of its table \verb|examinePrime c d p| (Definition~\ref{def:examinePrime}), for a prime
  power $p^e$ those of its tables \verb|examinePower c d p e| (Definition~\ref{def:examinePower}),
  and for a product those of the entry at its position. The store that the loop of
  \verb|sieving_info| makes from a store of the primes, with each entry made at its position, holds
  the rows of the references of every run whose entries begin with the entries made, and every entry
  that \verb|makeModulus?| makes at the position $n$ fits them.
\end{lemma}

\begin{proof}
  The coherence of the shared tables of rows is in $R$, a function of the reference: two factors
  with one reference have the same table of the prime or tables of the prime power, both
  \verb|examinePrime c d p| or both \verb|examinePower c d p e|, which the contract of the sieve
  entry of a modulus gives, so they have one closed form. A store of the primes holds no row but
  rows of $0$ of primes, which are \verb|sieves0| and the closed forms for every table of the prime
  (Lemma~\ref{lem:sieves0}); the tables of rows made with an entry are made only when absent and
  hold no row but the rows of $0$ of primes. No sizes of the tables of rows are needed: a row read
  or built is the closed form even where the table of rows is missing (Lemma~\ref{lem:fetchRow}).
  The fit of an entry needs only its modulus, its factors, its reference and its kind, so it
  survives the filling of the inverses of the strides $\mathrm{dinv}$ by the rows of offsets
  (Definition~\ref{def:classOffsets}).
\end{proof}

\begin{definition}[The missing rows of a denominator]\label{def:buildRows}
  \uses{def:fetchRow, def:rows}
  Proposed name: \lname{Ratpoints.Sieve.buildRows}. For the sieve specs of the first two stages that
  \verb|set_rows| makes (Definition~\ref{def:siftBits}), with the empty array for a row not built
  yet and the flag of a missing row: when a row is missing, for each spec whose row is not built
  yet, of the $n$-th sieve entry and its residue, the row of the residue from the table of rows of
  the entry, read again, and otherwise built into it by the builder of the entry
  (Definition~\ref{def:fetchRow}), with the end $p$; the specs and the store returned
  (\csource{find_points.c}{3522}{3554}).
\end{definition}

\begin{lemma}[The sieve specs of a denominator]\label{lem:buildRows}
  \uses{def:buildRows, lem:fetchRow, lem:storeTables, def:siftBits, lem:siftBits}
  Proposed names: \lname{Ratpoints.Sieve.BuildRowsPre}, \lname{Ratpoints.Sieve.BuildRowsInvariant},
  \lname{Ratpoints.Sieve.buildRows\_spec}, \lname{Ratpoints.Sieve.length\_setRows},
  \lname{Ratpoints.Sieve.getElem!\_setRows}, \lname{Ratpoints.Sieve.buildRows\_setRows\_spec},
  \lname{Ratpoints.Sieve.buildRows\_setRows\_map\_spec},
  \lname{Ratpoints.Sieve.RowStore.Fits.buildRows\_setRows},
  \lname{Ratpoints.Sieve.mapIdx\_eq\_map\_fst\_packedReads},
  \lname{Ratpoints.Sieve.eq\_map\_fst\_packedReads\_of\_eq\_run\_buildRows}. For a store that holds
  $R$ for the sieve entries of the first two stages, entries that fit $R$ at their positions and
  residues below their moduli: the sieve specs that \verb|set_rows| makes and \verb|buildRows|
  completes have the closed forms of the rows of their entries at their residues, the offsets of the
  class, the start $0$ and the end $p$, and the store after them still holds $R$; so a store that
  fits the entries fits them again after a denominator, and with the residues $b\,\mathrm{dinv}[k]
  \bmod p$ these specs are those whose readings Lemma~\ref{lem:siftBits} takes. This equation hides
  the store from the contracts below one denominator, which see sieve specs with the closed forms of
  their rows and do not mention the store.
\end{lemma}

\begin{proof}
  An invariant of the loop over the specs: the store holds $R$, and each spec has its closed form
  and the end $p$ or, from the position of the loop on, no row yet. A row present is $R$ at its
  reference, the closed form by the fit of the entry; a row built is the closed form
  (Lemma~\ref{lem:fetchRow}), and the store grows by rows of $R$, since the rows of the factors are
  $R$ at their references; the row of $0$ of a prime is in its table of rows, so a builder of a
  prime is called at a residue other than $0$ only.
\end{proof}

\verb|sift| takes the store and returns it with the rows built for the denominator: \verb|set_rows|
reads the rows from the tables of rows of the entries, \verb|buildRows| builds the missing ones, and
the bits of the interval are sifted with these specs (Definition~\ref{def:sift}). It takes the sieve
entries of the first two stages in place of their moduli with the thunks of the rows; its contract
takes, beside the entries, their reciprocals, their residues $b\,\mathrm{dinv}[k] \bmod p$, their
tables of inverses, and a store that fits them, as hypotheses on its arguments; it gives the store
back fitting them; the class word, the position of the row offsets of the class, and the packing of
the entries at those offsets remain hypotheses of the points equation, which is unchanged. The run
(Definition~\ref{def:sieveRun}) passes the store from one denominator to the next, the store fitting
the first $s_2$ entries in the invariant of its loop, and does not return it, nothing reading the
tables of rows after the run; the search (Definition~\ref{def:searchPoints}) hands the store of
\verb|sieving_info| to the run; and what the run needs of the search (Lemma~\ref{lem:soundRun})
gains the store with the field that it fits the entries of the first two stages, which the store of
\verb|sieving_info| gives, the entries changing after it only in the inverses of the strides.

\section{The command-line program and the differential tests}\label{sec:rowsCli}

The program does not change. The differential tests compare every point list in its order. Whole
chunks occur in the real runs only from heights of about $2000$ on, where the stride $1$ gives
pieces of $16$ bit arrays and more: at $H = 2000$ one whole chunk (a piece of $16$ bit arrays at the
stride $1$), at $5000$ two in the pieces of $40$ bit arrays of the stride $1$ and one in those of
$20$ of the stride $2$, at $30000$ up to $14$; so the suite of point-rich curves at each of its
heights $2000$, $5000$ and $30000$, the runs of \verb|test4.sh| at the height $30000$ and
\verb|test4high| run whole chunks. Below those heights every piece that reaches a first stage with a
modulus is a tail of fewer than $16$ bit arrays, so every suite that sieves runs tails, and the
suite of the sieving parameters runs the arm with no modulus of the first stage (\verb|-p 0|,
\verb|-n 0| and \verb|-n 0 -N 0|), which runs neither chunks nor tails. Pieces outside \verb|small|
occur only in the \verb|#guard| examples, since the tests stay far below its limit. Since
\verb|findPoints| now rounds the array size up to the chunk, its \verb|#guard| examples run each
class as one piece of $16$ bit arrays, and the examples with pieces of fewer than $16$ bit arrays
call \verb|sift| and the run directly.

The \verb|#guard| examples compare the chunked arm with the first stage one row at a time, word for
word with the place of the sentinel, on pseudo-random rows with their copies, for pieces of $1$ to
$40$ bit arrays (every tail, none, one and two whole chunks), moduli below the chunk and above it,
with \verb|small| and without, at word numbers far below $0$, and with no modulus; and the legs of
every tail of $t < 16$ bit arrays. They compare the row of $0$ with the pattern \verb|sieves0| of
the build of ratpoints for the primes $3$ to $13$, and over all its words with the words of the
forbidden divisor for every prime of the table; each builder with the closed form of its admission
test over all its words: every residue of the primes up to $13$ and four residues of the larger
primes up to $251$, with the tables of two curves and a table of all ones, every residue of the
prime powers $9$, $25$, $27$, $49$ and $121$ and six of $243$, and every residue of the products
$15$, $21$ and $45$ and some of $63$, $105$, $225$ and $255$, each builder storing its row into the
table of rows of its entry and no other row there, a product building the rows of its factors into
their tables of rows; the bits of \verb|fsq| and \verb|gsq| of the prime powers $9$ and $25$ for two
curves, read from a build of ratpoints, with none set from $m$ on; and the array size rounded up to
a multiple of the chunk, and the cap $2^{20}$, a multiple of it. They check the row of a sieve entry
read or built for the primes $3$ and $67$, the prime power $9$ and the products $15$ and $45$, which
share tables of rows with them; three denominators through \verb|sift| with one store, whose tables
of rows hold exactly the rows needed; and the tables of rows that \verb|sieving_info| makes, a prime
that says nothing among them.

\section{Deviations}\label{sec:rowsDeviations}

Each component differs from the C code in a stated way, recorded in its \verb|Deviation:| paragraph:
\begin{itemize}
  \item the chunk $16$, the value of the default build, set in the model, where the C code chooses
    it by the macros that the compiler is given, each build with vector registers taking $16$ and
    the build with $64$-bit bit arrays $1$ unless \verb|RATPOINTS_CHUNK| is given; the arm of the
    first stage chosen by a test on the constant, for the conditional compilation that compiles one
    arm only, and the guards of the legs of the tail on the chunk tested with those on $t$; no check
    at compile time that the chunk is at most $16$, the \verb|#error| of the C code, which its place
    inside the conditional compilation of the chunked arm makes dead code: a chunk above $16$
    compiles the other arm, and the model chooses the arm by the same test;
  \item the starts of the rows in a local array of the sieve specs, for the sieve specs of the
    sieving thread, whose starts the C code writes and nothing reads after the first stage; the
    reduction by the reciprocal under \verb|small| alone, for \verb|small && RP_USE_MOD_MUL|, with
    \verb|RP_USE_MOD_MUL| the constant $1$ of the production build (\csource{sift.c}{236}{238},
    excluded with the switch of the development builds); the sum $w_{\mathrm{low}} +
    \mathrm{offset}$ of the starts an integer that stays a \verb|long| where \verb|mod| reduces it
    outside \verb|small| (\verb|fitsLong_neg_sub_and_fitsLong_add|); one definition for the stepping
    back of a start at its two places (after the first modulus of a chunk and after each further
    one); the registers \verb|reg0|, \dots, \verb|reg15| a vector of $\mathrm{chunk}$ bit arrays,
    made of zeros and then loaded, and the loads, the ands and the stores of the registers loops
    over its bit arrays, for the variables and the code that the build unrolls under
    \verb|#if (RATPOINTS_CHUNK >= k)|; the pointers that each load and each and steps on by one bit
    array as the start plus the chunk; the array \verb|reg[8]| of a leg a vector of $8$ bit arrays,
    without the pragmas that make the compiler unroll its loops (\verb|RP_UNROLL_REGS|,
    \csource{sift.c}{466}{472}, excluded); indices into the rows and into the survivors for the
    pointers \verb|start|, \verb|end| and \verb|surv|, the chunk counted from the first bit array of
    the piece for the word number \verb|w_low_new|, and the end of the whole chunks an index for
    \verb|w_high_full|; no modulus tested by the size of the local array, for \verb|sp1 == 0|; the
    class word written into the bit arrays by a loop that counts up, bit array $n - 1 - i$ at the
    step $i$, for the loop that counts down from the last, in both arms; no output under
    \verb|DEBUG|; the other deviations of the first stage of Section~\ref{sec:packingDeviations}
    (the class word for \verb|RBA(packed)|, the moduli and the reciprocals as two lists) hold for
    both arms;
  \item one function for every prime below $64$ and one for every prime above $64$, with the prime
    that of its table, for the function that the build generates for each prime with the prime a
    constant, and the kind of the builder for the pointer \verb|init|, the call through it a match
    on the kind, and the kind set where the model makes the sieve entry of every kind of modulus
    (Definition~\ref{def:makeModulus}), where \verb|examine_prime| sets that of a prime when it
    makes its entry (\csource{find_points.c}{1800}{1800}), as the reciprocal, the bias and
    $\mathrm{rbainv}$ are; the table of the prime or the tables of the prime power of the factor and
    the reference to its entry as arguments, for the sieve entry \verb|se| from which the builders
    read them and into which they store the row, and the store of the tables of rows passed and
    returned, for the arrays \verb|sieve| of the entries; the row stored plainly, for the release
    store of \verb|RP_PUBLISH| in a build with threads, and read plainly by the builder of a
    product, for the acquire load \verb|RP_ACQUIRE| (stage~I); a fresh vector for each row, for the
    space that the builders take from the buffer of the tables (\verb|ba_next|), fresh arrays of
    zeros for \verb|help| and the pattern of a prime power, whose words the C code leaves unset
    until its loops clear them, and for \verb|rep|, of which \verb|lay_out_pattern| clears the first
    \verb|words| words and leaves the others unset and never reads them, the loops kept; the forms
    of the production builds, where \verb|RP_INIT_BRANCH| and \verb|RP_INIT_ONEWAY| are not set; the
    variables that the macros declare or read in their callers passed and returned; the steps
    \verb|ab += d; if (ab >= p) ab -= p;| and the like as conditional values; no copy \verb|pattern|
    of \verb|help0|; the variables \verb|ldp|, \verb|p1|, \verb|wp|, \verb|diff_shift| and
    \verb|diff| computed where they are used; the shifts by \verb|LONG_SHIFT| and the masks by
    \verb|LONG_MASK| as divisions and remainders by $64$; the loops that step by $m$ or by $q$ as
    \verb|while| loops; the number of the factors the length of their list, for \verb|nf|;
  \item the row of $0$ of a prime a function of $p$, computed when it is called, its first $(p +
    \mathrm{chunk} - 1)\,\mathrm{rbaPack}$ words, for the row of the generated table, whose other
    words are $0$ and never read; the kind of the builder of a prime a function of $p$, for the
    generated table of pointers; the words of a forbidden divisor one function for the array
    \verb|work| of the generator and for the words that \verb|sieving_info| makes for the primes
    beyond the table, with the positions $p\,i$ for $i < 64$, for the positions $m = 0, p, 2p,
    \dots$ below $64p$ (line 2776), in a fresh array of $p$ words, for \verb|work| of
    \verb|RATPOINTS_MAX_PRIME| words and the space in \verb|forb_words|, and those of a prime of the
    table an array of the first $p$ words of the row of $0$, for the pointers \verb|start| and
    \verb|end| into it;
  \item the tables of rows in a store that maps the reference to an entry to its table of rows, for
    the arrays \verb|sieve| of the entries, which the C code reaches through pointers
    (\verb|sieve_list[n]| and the field \verb|factor| of a product), none for \verb|NULL|: a prime
    by itself for its index in the table, a prime power by its prime and its exponent for the index
    of its entry in \verb|pw_idx[pn][e]|, a product by its position among the moduli; the table of
    rows of a prime made by the loop of the tables when the prime says something, the test of line
    1764 made there on the table of the prime; the tables of rows of a prime power and of a product
    made by the loop over the moduli right after the entry, where \verb|make_power| and
    \verb|make_modulus| make them with the entry, each only when its reference is absent, as the
    entry of a prime power is made once; a table of rows made also for a prime beyond those of the
    loop of the tables and for one that says nothing, for which \verb|examine_prime| makes no entry
    and which ratpoints never takes as a factor; the store not returned by the run, which nothing
    reads after it;
  \item in \verb|set_rows|, the empty array for the null pointer of a row not built yet; the start
    $0$, where the C code leaves the field as it is until the first stage sets it; the end $p$ also
    for a row not built yet, where the C code adds $p$ to the null pointer; the row read plainly,
    for the acquire load that the flag \verb|shared| chooses, and one call, for the two calls with
    the flag a constant (stage~I); in the loop that builds the missing rows, no lock of the tables
    (\verb|rp_tables_lock|, \verb|rp_tables_unlock|) and the pointer read again plainly, for the
    acquire load, both for the threads that share the tables of rows (stage~I), the reading again
    itself modelled, since with one thread it finds a row that an earlier spec of the same
    denominator built into the same table of rows: when a prime or a prime power occurs twice among
    the moduli of the first two stages, which then have one table of rows, or when it comes after a
    product among them that has it as a factor, whose builder builds the row of the factor into the
    table of rows of the factor (cases that the parameters admit and ratpoints never has, since it
    never takes two moduli with a common prime); a product that occurs twice has a table of rows of
    its own at each position, so that its reading again finds nothing; no measurement of the
    construction (\verb|RP_INIT_TIC|, \verb|RP_INIT_TOC|) and no output under \verb|DEBUG|;
  \item the bit arrays \verb|sq|, \verb|fsq| and \verb|gsq| of \verb|examine_power| fresh arrays of
    $\mathrm{modWords}$ words, with $64$ for \verb|LONG_LENGTH| in $\mathrm{modWords}$, which the
    loop at line 2282 clears, for the local array and the arrays of the entry \verb|pw|; the
    variable $t$ for the multiples of $p$, where the C code reuses $x$; the bits of the squares and
    the inverses functions of their own, the bit array passed and returned and the inverses a fresh
    array returned; the deviation of sub-stage~F1 for the Booleans of \verb|fsq| and \verb|gsq|
    goes;
  \item natural numbers, and integers for the starts, for the \verb|long|s of these components, the
    product $b\,\mathrm{inv}[x]$ of the builder of a prime power below $m^2$; the range obligations
    of \verb|examine_power| keep their lemmas (Section~\ref{sec:tablesDeviations}).
\end{itemize}
The lines of this part of the C code that are neither modelled nor left to a later stage are
excluded in \verb|correspondence/ratpoints.toml|: the pragma \verb|RP_UNROLL_REGS| of
\verb|tail_leg| (\csource{sift.c}{466}{472}), which carries no meaning; the alignment of the buffers
and the sizing, allocation and enlargement of the buffer of the rows of the tables, the number of
the prime powers that sizes the buffers, and the clearing of the pointers to them
(\csource{find_points.c}{181}{270}); the function \verb|find_points_init|
(\csource{find_points.c}{272}{375}), which a caller of the library calls before
\verb|find_points_work| to allocate the arrays of the arguments, and whose refusal of a degree
\verb|find_points_work| repeats (lines 4557-4558, which the model has; the reason of the exclusion
of lines 4618-4623 says the same); and \verb|find_points_clear| (\csource{find_points.c}{377}{416}),
which frees them; the model allocates nothing. Left to stage~I are the lock of the tables and the
pool of threads (among them the declarations at \csource{find_points.c}{41}{43}), the orderings of
the acquire loads and the release stores, and the cost model of the choice of the moduli, which
reads the array size rounded to the chunk, so that ratpoints chooses its moduli by the chunk too:
the moduli of the model are those of the parameters.

\section{What changes in the earlier chapters}\label{sec:rowsEarlier}

\begin{itemize}
  \item Chapter~\ref{chap:sieve}: the rows of a denominator (Definition~\ref{def:admits}) lose their
    links to \verb|init.c|, \verb|gen_init_sieve_h.c| and \verb|gen_find_points_h.c|, which go to
    the builders and the row of $0$ of Chapter~\ref{chap:rows}: the admission test, with no C range,
    is what the rows are specified by through the closed form (Definition~\ref{def:closedRow}); the
    paragraph on the rules of the translation no longer leaves the generated tables to stage~F4, the
    row of $0$ being proved the closed form (Lemma~\ref{lem:sieves0}); the run
    (Definition~\ref{def:sieveRun}) passes the store of the tables of rows from one denominator to
    the next, its contract and the reduction of the run to the brute-force reference
    (Theorem~\ref{thm:run_reference}) taking a store that fits the first $s_2$ entries, its node
    using Lemma~\ref{lem:buildRows}; the introduction counts the chunked first stage and the rows of
    the builders among the data refinements; the record of comparator after
    Theorem~\ref{thm:findPoints}, with the files of the rows among those without Mathlib; the files
    section names \verb|Ratpoints/Model/Rows.lean| for the rows.
  \item Chapter~\ref{chap:words}: the introduction describes the default build with the chunk $16$
    and both arms in place of the build \verb|chunk1|, and its paragraph on \verb|array_size| names
    \verb|find_points_work_1| for the caps and the rounding; the pieces
    (Definition~\ref{def:arraySize}) are capped at $\lfloor 2^{20}/\mathrm{chunk} \rfloor \cdot
    \mathrm{chunk}$ and rounded up to a multiple of the chunk,
    \verb|Ratpoints.Params.one_le_arraySizeAt| using $0 < \mathrm{chunk}$, the node using
    Definition~\ref{def:chunk}, and the paragraph after it gives the chunk its value; the row of a
    denominator in words (Definition~\ref{def:row}) is the closed form of $p + \mathrm{chunk} - 1$
    bit arrays that the builders build, with \verb|Ratpoints.Sieve.SieveEntry.admits_congr|,
    \verb|admits_add_mul| and \verb|admits_mod| and no C range, its node using
    Definition~\ref{def:SieveEntry} only, no longer Definition~\ref{def:rbaPack}, and naming
    Definition~\ref{def:closedRow} for the closed form without using it: the plain builder
    (\verb|packedPattern|, \verb|SieveEntry.row|, \verb|SieveEntry.getLsbD_row|) retires, and
    \verb|window| with \verb|getLsbD_window| goes to Definition~\ref{def:layOutPattern} and
    Lemma~\ref{lem:compositeRows}; the rows of a modulus (Definition~\ref{def:rows}) keep
    \verb|SieveSpec| with the start and the end, in \verb|Ratpoints/Model/FirstStage.lean|, and
    \verb|SieveSpec.wordAt|, the node using Definition~\ref{def:rbaPack} in place of
    Definition~\ref{def:row}, the thunks of the rows (\verb|SieveEntry.rows|,
    \verb|SieveEntry.get_rows|, \verb|SieveEntry.row_mod|) retiring for the tables of rows of
    Definition~\ref{def:entryRef}; the operations on bit arrays (Definition~\ref{def:bitArrayOps})
    say that \verb|RBA| goes with both arms; the first stage (Definition~\ref{def:andRows}) writes
    the class word from the last bit array (\csource{sift.c}{954}{954}, within its range) and is the
    arm of a chunk of $1$ or above $16$, its node using Definition~\ref{def:bitArrayMoves}; one
    piece (Definition~\ref{def:siftPiece}) chooses the arm, its node using
    Definition~\ref{def:andRowsChunked} and Lemma~\ref{lem:andRowsChunked}; one denominator
    (Definition~\ref{def:sift}) takes the sieve entries of the first two stages in place of their
    moduli with the thunks of the rows, and takes and returns the store, the node using
    Definition~\ref{def:buildRows} and Lemma~\ref{lem:buildRows}, gives
    \csource{find_points.c}{3522}{3554} to Definition~\ref{def:buildRows} and keeps
    \csource{find_points.c}{3490}{3521} and \csource{find_points.c}{3555}{3561}; its contract and
    Lemma~\ref{lem:sift_eq_checkDen} take those sieve entries, no longer the rows of sieve entries,
    and a store that fits them, the contract giving the store back fitting them, the points equation
    unchanged; the paragraph after that lemma speaks of the store; the list of deviations loses the
    items on the pattern written from the first bit array, on the first stage of one register per
    chunk, on the rows from the admission test and on the rows made when first needed; the files
    section names \verb|Ratpoints/Model/FirstStage.lean| and \verb|Ratpoints/Proof/FirstStage.lean|.
  \item Chapter~\ref{chap:setup}: the table of a prime with the exit
    (Definition~\ref{def:examinePrimeExit}) gives \csource{find_points.c}{1820}{1821} to
    Definition~\ref{def:storeTables} and keeps \csource{find_points.c}{1795}{1819} and
    \csource{find_points.c}{1822}{1826}; the search (Definition~\ref{def:searchPoints}) hands the
    store of \verb|sieving_info| to the run, the node using Lemma~\ref{lem:storeTables}; the
    paragraph after the proof of Theorem~\ref{thm:refusals}, on Theorem~\ref{thm:run_reference},
    adds a store of the tables of rows that fits the first $s_2$ entries to its list of hypotheses;
    the record of \verb|find_points_work_1| lacks only what belongs to stages~H and~I, with the
    exclusions of Section~\ref{sec:rowsDeviations}.
  \item Chapter~\ref{chap:denominators}: the forbidden divisors (Definition~\ref{def:forbidden})
    give the words of a forbidden divisor (\verb|forbiddenRow|) with
    \csource{find_points.c}{2773}{2777} and their bit lemma to Definition~\ref{def:sieves0} and
    Lemma~\ref{lem:sieves0}, which they use, and their link to \verb|gen_find_points_h.c| to
    Definition~\ref{def:sieves0}; the words of a prime of the table are the first $p$ words of the
    row of $0$, their lemmas going through Lemma~\ref{lem:sieves0}; the paragraph on the run before
    Section~\ref{sec:denSearch} says that the run also takes the store of the tables of rows and
    passes it from one denominator to the next, the rows of the first $s_2$ entries read from it or
    built, and its contract (\verb|sieveRun_spec|, \verb|toList_run_sieveRun|,
    \verb|toList_run_sieveRun_eq_flatMap|) adds a store that fits the first $s_2$ entries; the
    search paragraph of Section~\ref{sec:denSearch} says that \verb|sieving_info| returns the store
    and the search hands it to the run; the files section puts the words in
    \verb|Ratpoints/Model/Rows.lean| and their bit lemma in \verb|Ratpoints/Proof/Rows.lean|.
  \item Chapter~\ref{chap:tables}: the introduction refers to Chapter~\ref{chap:rows} for F4; the
    tables of a prime power (Definition~\ref{def:examinePower}) keep \verb|fsq|, \verb|gsq| and
    \verb|sq| as bit arrays of words, the node using Definition~\ref{def:powerWords}, with
    \csource{find_points.c}{2267}{2282}, \csource{find_points.c}{2285}{2301} and
    \csource{rp-private.h}{531}{534}, giving \csource{rp-private.h}{523}{523},
    \csource{find_points.c}{2283}{2284} and \csource{find_points.c}{2302}{2303} to
    Definition~\ref{def:powerWords} and \csource{init.c}{359}{386} and
    \csource{rp-private.h}{592}{592} to Definition~\ref{def:sieveInitPower}; its lemma
    (Lemma~\ref{lem:power_admits}) names \verb|bitAt_fsq_run_examinePower_of_isSquare| and
    \verb|bitAt_gsq_run_examinePower_of_isSquare| in place of the lemmas on Booleans and loses
    \verb|ExaminePowerInvariant|, its contract and the lemma on the points taking $p^e \le
    64\,\mathrm{modWords}$, its proof using Lemma~\ref{lem:powerWords}; the sieve entries
    (Definition~\ref{def:SieveEntry}) gain the new fields \verb|ref| and \verb|init|, the reference
    to their table of rows and the kind of their builder, the node gaining
    \verb|Ratpoints.Sieve.Factor.ref|, \verb|Ratpoints.Sieve.Factor.init|,
    \verb|Ratpoints.Sieve.SieveEntry.refInit| and \verb|Ratpoints.Sieve.Factor.modulus|, the modulus
    of a factor, now a definition of the model, which moves here from
    Lemma~\ref{lem:sieveEntry_admits}: the reference of a factor, to which the field \verb|factor|
    of the entry of a product points, is that of its prime or of its prime power; the kind of the
    builder of the entry of a factor is, for a prime, the kind of the builder of the prime of its
    table, which \verb|examine_prime| sets (line 1800), and for a prime power the builder of a prime
    power, which \verb|make_power| sets (line 2420); and the reference and the kind of the builder
    of the modulus at the position $n$ with its factors are, for one factor (a prime or a prime
    power), those of the entry of that factor, whose entry the modulus shares, as
    \verb|sieve_list[n]| is \verb|prime_se[pn]| or \verb|power_se[idx]|, and for a product the
    reference to its own entry, by the position $n$, and the builder of a product; the node uses
    Definitions~\ref{def:entryRef} and~\ref{def:rowInit} and gives \csource{rp-private.h}{497}{498}
    to Definition~\ref{def:rowInit} and \csource{init.c}{399}{414}, \csource{init.c}{419}{429} and
    \csource{rp-private.h}{593}{593} to Definition~\ref{def:sieveInitProduct}; the sieve entry of a
    modulus (Definition~\ref{def:makeModulus}) takes the position and sets the reference and the
    kind of builder of its position (Definition~\ref{def:SieveEntry}), with
    \csource{find_points.c}{2414}{2426}, \csource{find_points.c}{2428}{2429},
    \csource{find_points.c}{2434}{2470} and \csource{find_points.c}{2472}{2475}, giving lines 2427
    and 2471 to Definition~\ref{def:storeTables}; its contract (Lemma~\ref{lem:sieveEntry_admits})
    gives the reference and the kind of builder of the position, the node using
    Lemma~\ref{lem:fetchRow} for the buildable factors, and \verb|Ratpoints.IsModulusEntry| and
    \verb|Ratpoints.Sieve.Factor.IsSound| gain that the factors are buildable, that their tables are
    those of \verb|examinePrime| and \verb|examinePower|, and that the moduli of the factors
    multiply to the modulus; the paragraph after that lemma has the builders of
    Chapter~\ref{chap:rows} in place of the plain builder; the loop of the tables
    (Definition~\ref{def:examinePrimes}) returns the store with the tables of rows of its primes,
    the node using Definition~\ref{def:storeTables}, and its lemma (Lemma~\ref{lem:examinePrimes})
    gains \verb|Ratpoints.Sieve.primeStore_of_run_examinePrimes_eq_some| next to its other lemmas on
    what the loop gives, with the statement that the store of the loop is a store of the primes
    (\verb|PrimeStore|: it holds every $R$ whose rows of $0$ of the primes are \verb|sieves0|, and
    every table of rows of a prime in it has its row of $0$), the node using
    Lemma~\ref{lem:storeTables}; \verb|sieving_info| (Definition~\ref{def:sievingInfo}) makes the
    tables of rows of each entry and returns the store, the node using
    Definition~\ref{def:storeTables}; its contract (Lemma~\ref{lem:sievingInfo}) takes the result
    with the store, and its postcondition becomes the structure \verb|Ratpoints.IsSievingInfo|, the
    six former conjuncts and the new field \verb|fits|, that the store fits the entries, the loop
    invariant \verb|Ratpoints.SievingInfoInvariant| holding, before a return, the structure
    \verb|Ratpoints.SievingInfoEntries| with the store made for the entries, so that the store that
    \verb|sieving_info| returns fits its entries; the node gains \verb|Ratpoints.IsSievingInfo|,
    \verb|Ratpoints.SievingInfoEntries|, \verb|Ratpoints.SievingInfoEntries.fits| and
    \verb|Ratpoints.refInit_push| and uses Lemma~\ref{lem:storeTables}; the deviations on the rows
    of a prime power and of a product and on the entry of a prime power kept once say what F4 does.
  \item Chapter~\ref{chap:stages}: the second stage (Definition~\ref{def:secondStage}) uses
    Definition~\ref{def:useCLong}, which it reached through Definition~\ref{def:row} before; the
    paragraph before Lemma~\ref{lem:staged} gives Lemma~\ref{lem:sift_eq_checkDen} a store that fits
    the entries; the paragraph on the run before Section~\ref{sec:cuts} says that the run reads or
    builds the rows of the first $s_2$ entries through the store, adds the fitting store to the
    hypotheses of its contract, and says that the search passes the store of \verb|sieving_info|,
    which Lemma~\ref{lem:sievingInfo} gives fitting the entries; the deviations no longer leave the
    chunked first stage and the ends of the rows to stage~F4.
  \item Chapter~\ref{chap:packing}: the paragraph on \verb|_ratpoints_sift0| in the section on the C
    code says that the start $r = (-w_{\mathrm{low}} - \mathrm{offset}) \bmod p$ at bit array $p -
    r$ is that of the arm one row at a time (Definition~\ref{def:andRows}), the arm of a chunk of
    $1$ or above $16$, and that the chunked arm of the default build, the chunk $16$ of the model,
    starts each modulus at $(w_{\mathrm{low}} + \mathrm{offset}) \bmod p$, by \verb|mod_mul| under
    \verb|small| and by \verb|mod| otherwise (\csource{sift.c}{596}{606},
    Definition~\ref{def:chunkStarts}); the rows read in the packing (Lemma~\ref{lem:reads_row})
    speak of sieve specs of five fields with the closed form of the row, the node using
    Definition~\ref{def:row} and, in its proof, Lemma~\ref{lem:closedRow}; a sieve spec that fits
    (Lemma~\ref{lem:rowIndices}) has the end $p$ and a row that wraps, the node using
    Definition~\ref{def:closedRow} through the constructor of \verb|SieveSpec.Fits|, whose field
    says that the row wraps, and the lemma gains
    \verb|SieveSpec.Fits.ite_modMul_modLong_add_eq_emod|: the start of the chunked arm, by the same
    reductions, is $(w_0 + \mathrm{offset}) \bmod p$, the sentence on $p - r$ naming the arm one row
    at a time; the paragraph after that lemma has either arm, Definition~\ref{def:andRows} or
    Definition~\ref{def:andRowsChunked}, chosen by one piece, as the first stage, both taking the
    offsets and the reciprocals from the position $0$ and their contracts giving the same words; the
    bits of one denominator (Definition~\ref{def:siftBits}) have \verb|set_rows| read the rows from
    the store, with the empty array for a row not built yet, the start $0$, the end $p$ and the flag
    of a missing row, the node using Definition~\ref{def:entryRef}; its lemma
    (Lemma~\ref{lem:siftBits}) builds the specs with the closed forms, \verb|setRows_map| giving way
    to the lemmas of Lemma~\ref{lem:buildRows}; the paragraph before Lemma~\ref{lem:soundRun} adds
    the fitting store to the hypotheses of the contracts of one denominator and of the run, and says
    that the search hands on the store that \verb|sieving_info| returns; what the run needs of the
    search (Lemma~\ref{lem:soundRun}) gains the store, the node using Lemma~\ref{lem:storeTables};
    the deviation that writes out $-w_{\mathrm{low}} - \mathrm{offset}$ in both reductions of the
    first stage names the arm one row at a time, the \verb|long| $a = w_{\mathrm{low}} +
    \mathrm{offset}$ of the chunked arm being covered by the deviation of
    Definition~\ref{def:chunkStarts}; the deviations and the work left to stage~F4 go; the files
    section counts the lines anew.
  \item The stage list of the introduction refers to this chapter (Chapter~\ref{chap:rows}), which
    is input after Chapter~\ref{chap:packing}. The files sections, the uses and the generated table
    follow.
\end{itemize}

\section{Files, names and order of work}\label{sec:rowsFiles}

The model, in the folder \verb|Ratpoints/Model|: the chunk goes into \verb|Constants.lean| with
$\mathrm{modWords}$, each citing its macro; the bits of a bit array of words, the squares and the
inverses of a prime power, the references, the kinds of builders with that of a prime for the
generated table \verb|sieve_init|, and the modulus of a factor into \verb|Tables.lean|; the sieve
specs, the and of two bit arrays with the load, the store and the fill, and both arms of the first
stage into a new file \verb|FirstStage.lean| below \verb|Sieve.lean| (about $470$ lines); the row of
$0$, the words of a forbidden divisor, the store, the builders and the reading and building of the
row of a sieve entry into a new file \verb|Rows.lean| above it (about $870$ lines); the missing rows
of a denominator stay with \verb|sift| in \verb|Sieve.lean|, with the clearing of the low and the
high bits and the test for zero (Definition~\ref{def:bitArrayOps}), which shrinks to about $1000$
lines. The proofs, in the folder \verb|Ratpoints/Proof| and all without Mathlib but for those of the
entries of the moduli, the loop of the tables, the run and the search: the facts on vectors, arrays,
lists, ranges and the natural numbers and integers that index them in a new core file
\verb|Containers.lean|, split off \verb|Bits.lean| and below it, \verb|Bits.lean| keeping the facts
on \verb|UInt64| and \verb|Int64| (about $1100$ and $420$ lines); the closed form, the row of $0$
and the bit lemma of the words of a forbidden divisor, which leaves \verb|Forbidden.lean|, in a new
file \verb|Rows.lean| below \verb|Words.lean|; the end and the copies of a fitting spec in
\verb|Reductions.lean|; the loads, the stores and the fill of the first stage in \verb|Words.lean|;
the chunked arm in a new file \verb|FirstStage.lean| above it; the builders of a prime in a new file
\verb|PrimeRows.lean| and those of a prime power and of a product in a new file
\verb|CompositeRows.lean| above it; the store, the builders through it, its invariant and the
missing rows of a denominator in a new file \verb|RowStore.lean| above them; the bit arrays of a
prime power and the entry of a modulus with its reference in \verb|Moduli.lean|; the store of the
loop of the tables in \verb|Forbidden.lean|; the arms of one piece and one denominator on the store
in \verb|Sift.lean|; the run in \verb|Denominators.lean|; and the store of \verb|sieving_info| and
what the run needs of the search in \verb|FindPoints.lean|. The general facts on core types go to
\verb|Containers.lean| and \verb|Bits.lean| under their names in the namespaces of Lean core: among
the helpers of the prototype, \verb|ite_add_mod_mod| and \verb|ringNext_eq| as lemmas on natural
numbers, \verb|all_range_and|, \verb|all_range_congr| and \verb|all_range_length_getElem!|,
generalized from the factors to any list, as lemmas on lists (once it is checked that core has none
of them), \verb|getElem!_append| and the entries of a fold of appends as lemmas on arrays and
\verb|toNat_emod_add_mod| as a lemma on integers; a helper that one lemma gives is inlined at its
uses instead (\verb|sub_mod_chunk_mod| by \verb|Nat.sub_mod_eq_zero_of_mod_eq|, \verb|add_le_of_dvd|
by \verb|Nat.le_of_lt_add_of_dvd|, \verb|ite_add_eq_mod| and \verb|add_mod_ring| by
\verb|Nat.mod_eq_ite_of_lt_two_mul|); the facts on the constants of the model and on the four lanes
of the builders of a prime stay in their files. Every file stays below $1500$ lines (the largest new
one, \verb|RowStore.lean|, about $1350$). The eight new modules join their part roots and
\verb|Ratpoints.lean|. The docstrings say ``table of rows'' for a table of the store, keep ``the
table of a prime'' and ``the tables of a prime power'' for the data of
Definitions~\ref{def:examinePrime} and~\ref{def:examinePower}, and use ``reference'' only for a
reference to an entry. One commit per item, each passing the project's checks, the build with other
values of the constants and the differential tests, and leaving the output unchanged, with the nodes
of the earlier chapters that it changes and the C links moved from their old owners in the same
commit:
\begin{enumerate}
  \item[F4.1] The chunk with $0 < \mathrm{chunk}$ (Definition~\ref{def:chunk}); the largest array
    size and the array size rounded to the chunk, with the lower bound of the array size; the
    \verb|#guard| examples at the chunk $16$, among them the example of \verb|findPoints| with
    pieces of $1$, $2$ and $3$ bit arrays restated: one array size in place of the three, which all
    round to $16$, at the same heights, its comment saying that each class runs as one piece of $16$
    bit arrays with the low and the high mask on the same piece, the last bit array cut at
    $w_{\mathrm{high}}$ inside the piece and no low mask when $H$ is a multiple of
    $\mathrm{rbaLength}$, and that several pieces and masks on different pieces are tested at the
    level of \verb|sift| (the example with pieces of one, two and three bit arrays at $H = 300$), by
    the example of the run with the array size $1$, and through \verb|findPoints| by the
    differential tests at the heights $5000$ and $30000$ and \verb|test4high|; the build with other
    values of the constants varies the chunk ($1$ and $17$ with the other constants, $2$, $3$, $5$
    and $12$ alone) and the width of the bit arrays alone ($\mathrm{rbaPack}$ $8$ and $2$, and $1$
    with the chunk $5$), and no longer the largest array size, which is derived; the node of
    Definition~\ref{def:arraySize}. About $100$ lines.
  \item[F4.2] The split of \verb|Ratpoints/Proof/Bits.lean| before it reaches $1500$ lines: the
    facts on vectors, arrays, lists, ranges, natural numbers and integers into the new core file
    \verb|Ratpoints/Proof/Containers.lean| below it, the facts on \verb|UInt64| and \verb|Int64|
    staying, both in the namespaces of Lean core, with the module docstrings, the row of the layout
    table of \verb|README.md|, the part root and \verb|Ratpoints.lean|; the type dump unchanged.
    About $1100$ lines moved.
  \item[F4.3] The sieve specs with the start and the end, the and of two bit arrays, the fill of a
    bit array (Definition~\ref{def:bitArrayMoves}, begun) and the first stage one row at a time,
    writing from the last bit array, in \verb|Ratpoints/Model/FirstStage.lean|; \verb|set_rows| with
    the start $0$ and the end $p$; the plain row with its copies, $p + \mathrm{chunk} - 1$ bit
    arrays; the closed form (Definition~\ref{def:closedRow}, Lemma~\ref{lem:closedRow}) in
    \verb|Ratpoints/Proof/Rows.lean|, with a lemma, until F4.10, that the plain row is the closed
    form, in the node of Definition~\ref{def:row} in place of \verb|SieveEntry.getLsbD_row| there
    and of \verb|SieveEntry.row_mod| in that of Definition~\ref{def:rows}; a fitting spec with its
    end and its copies, and the rows read in the packing, the first stage and the specs of a
    denominator on the five fields and the closed forms; the nodes of Definitions~\ref{def:row},
    \ref{def:rows} and~\ref{def:andRows} and Lemmas~\ref{lem:reads_row}, \ref{lem:rowIndices}
    and~\ref{lem:siftBits}. About $900$ lines.
  \item[F4.4] The chunked first stage (Definitions~\ref{def:chunkStarts}, \ref{def:chunkSieve},
    \ref{def:tailLegs} and~\ref{def:andRowsChunked}, Lemma~\ref{lem:andRowsChunked}), with the load
    and the store of a bit array (Definition~\ref{def:bitArrayMoves}),
    \verb|Ratpoints/Proof/FirstStage.lean|, and the choice of the arm in one piece; the exclusion of
    \verb|RP_UNROLL_REGS|; the first stage one row at a time the arm of a chunk of $1$ or above
    $16$, its \verb|Deviation:| paragraph losing the clause on the first stage of the builds with
    one register per chunk, and the generated table of the correspondence refreshed; the
    differential tests at the heights with whole chunks; comparator; the nodes of
    Definitions~\ref{def:siftPiece}, \ref{def:bitArrayOps} (\verb|RBA| with both arms)
    and~\ref{def:andRows} (the arm of a chunk of $1$ or above $16$) and of
    Lemma~\ref{lem:rowIndices} (the start of the chunked arm). About $1450$ lines.
  \item[F4.5] The tables of a prime power as words (Definition~\ref{def:powerWords},
    Lemma~\ref{lem:powerWords}), with $\mathrm{modWords}$, \verb|setSquareBits| and
    \verb|unitInverses|, which take \csource{rp-private.h}{523}{523},
    \csource{find_points.c}{2283}{2284} and \csource{find_points.c}{2302}{2303} from the tables of a
    prime power, and \verb|Ratpoints.Sieve.Factor.modulus| moved into the model, claimed by
    Definition~\ref{def:SieveEntry}; the nodes of Definitions~\ref{def:examinePower}
    and~\ref{def:SieveEntry} and Lemmas~\ref{lem:power_admits} and~\ref{lem:sieveEntry_admits}.
    About $500$ lines.
  \item[F4.6] The words of a forbidden divisor, the row of $0$ and the builders of a prime
    (Definitions~\ref{def:sieves0}, \ref{def:rowInit}, \ref{def:initMacros}, \ref{def:sieveInit1}
    and~\ref{def:sieveInit2}, Lemmas~\ref{lem:sieves0} and~\ref{lem:primeRows}), with the
    references, the store and its lookup and put (Definition~\ref{def:entryRef}),
    \verb|Ratpoints/Model/Rows.lean| with the words of a forbidden divisor and \verb|window| moved
    in, their bit lemma moved to \verb|Ratpoints/Proof/Rows.lean|,
    \verb|Ratpoints/Proof/PrimeRows.lean| and the beginning of \verb|Ratpoints/Proof/RowStore.lean|:
    the run does not call the builders yet, and \verb|#guard| examples compare them with the plain
    row, which, with the admission test of a prime, gives its C lines of the primes to the builders;
    the nodes of Definitions~\ref{def:admits}, \ref{def:row}, \ref{def:forbidden}
    and~\ref{def:SieveEntry}. About $1500$ lines.
  \item[F4.7] The builders of a prime power and of a product and the row of a sieve entry read or
    built (Definitions~\ref{def:layOutPattern}, \ref{def:sieveInitPower}, \ref{def:sieveInitProduct}
    and~\ref{def:fetchRow}, Lemmas~\ref{lem:compositeRows} and~\ref{lem:fetchRow}), with the
    references and the kinds of builders of the factors and of the moduli and the fields \verb|ref|
    and \verb|init| of the sieve entries in Definition~\ref{def:SieveEntry}, set by the sieve entry
    of a modulus at its position, \verb|Ratpoints/Proof/CompositeRows.lean| and the builders through
    the store in \verb|Ratpoints/Proof/RowStore.lean|; still not called by the run, and compared
    with the plain row, which, with the admission tests of a prime power and of a product, gives
    them its last C lines; the nodes of Definitions~\ref{def:row}, \ref{def:examinePower},
    \ref{def:SieveEntry} and~\ref{def:makeModulus} and Lemma~\ref{lem:sieveEntry_admits}. About
    $1800$ lines.
  \item[F4.8] The forbidden divisors of the primes of the table read the row of $0$; the node of
    Definition~\ref{def:forbidden}. About $60$ lines.
  \item[F4.9] The tables of rows made with the entries (Definition~\ref{def:storeTables},
    Lemma~\ref{lem:storeTables}): the loop of the tables and \verb|sieving_info| make the tables of
    rows and return the store, which the search does not hand on yet; the invariant of the store,
    the store of the loop of the tables and the fit of the store that \verb|sieving_info| returns,
    with the soundness of the entries of the moduli extended; the \verb|Deviation:| paragraph of
    \verb|examinePrime?| saying that the table of rows of a prime that a modulus needs (one beyond
    the loop of the tables, which \verb|makeModulus?| examines, or one that says nothing) is made by
    \verb|SieveEntry.addTables| in the loop of \verb|sieving_info| over the moduli, right after the
    entry; the exclusions of the buffers (\verb|find_points.c:181-270|, memory), of
    \verb|find_points_init| (\verb|find_points.c:272-375|, interface) and of
    \verb|find_points_clear| (\verb|find_points.c:377-416|, memory), with the reason of the
    exclusion of \verb|find_points.c:4618-4623| restated: \verb|find_points_work| refuses an
    inadmissible degree first, at lines 4557-4558, so that a missing work array there comes only
    from a caller who changed the degree after \verb|find_points_init|; the nodes of
    Definitions~\ref{def:examinePrimeExit}, \ref{def:examinePrimes}, \ref{def:makeModulus}
    and~\ref{def:sievingInfo} and Lemmas~\ref{lem:examinePrimes}, \ref{lem:sieveEntry_admits}
    and~\ref{lem:sievingInfo}. About $1200$ lines.
  \item[F4.10] The run on the store: \verb|set_rows| reads the rows from the tables of rows, the
    missing rows of a denominator are built (Definition~\ref{def:buildRows},
    Lemma~\ref{lem:buildRows}), \verb|sift| and the run pass the store on and the search hands it to
    the run; the plain builder and the thunks of the rows retire, with the lemma of F4.3; the
    contracts of one denominator, of the run and what the run needs of the search on the store;
    comparator; the nodes of Definitions~\ref{def:sieveRun}, \ref{def:row}, \ref{def:rows},
    \ref{def:sift}, \ref{def:searchPoints} and~\ref{def:siftBits}, Theorem~\ref{thm:run_reference}
    and Lemmas~\ref{lem:sift_eq_checkDen}, \ref{lem:siftBits} and~\ref{lem:soundRun}. About $1000$
    lines.
  \item[F4.11] The passages of the earlier chapters that F4 makes untrue
    (Section~\ref{sec:rowsEarlier}) beyond the nodes of the items; the records of comparator; the
    description of the repository and \verb|formalization.yaml|.
\end{enumerate}
F4.9 and F4.10 together are the switch to the store; the cut between them changes each contract
once, that of \verb|sieving_info| in F4.9 and those of one denominator and of the run in F4.10, and
leaves for one commit a store that the search makes and does not use.

\section{The evidence}\label{sec:rows-evidence}

\begin{itemize}
  \item A reading of the C code of the chunk, the first stage and the rows at the pin, with a second
    reader who checked it, on instrumented builds of the C program: the survivors of the first stage
    equal, word for word, those of the closed forms, on the runs of \verb|test4.sh| and $900$ random
    runs in the default build with the chunk $16$, the build with the chunk $1$ and one with the
    chunk $12$, and on $36$ runs at large heights in the first two; in each build about $3$ million
    pieces and $306$ million bit arrays, the arm with no modulus of the first stage, every length of
    the tail from $0$ to $15$ and pieces shorter than a chunk among them, and $22$ pieces outside
    \verb|small| in \verb|test4.sh|, which take the path through \verb|mod|; about $235\,000$ rows
    read during these runs, each with all its $m + \mathrm{chunk} - 1$ bit arrays, \verb|sieves0|
    and the rows of primes, prime powers and products, equal their closed forms; a build with the
    chunk $17$ compiles the other arm and gives the output of \verb|testbase4|, the \verb|#error|
    never firing. The reading finds one more defect of the C code without effect on the points, now
    recorded with ratpoints: \verb|set_rows| adds $p$ to the null pointer of a row not built yet
    (line 3166), which the C standard leaves undefined and gcc compiles harmlessly.
  \item Before the sub-stage, a sketch of the model in Lean, in three parts (the chunked first
    stage; the builders and the tables of a prime power as words; the store and the switch), which
    agrees with \verb|ratpoints -q| on the differential tests, every point list in its order, $3315$
    comparisons in the default suites and $13$ in \verb|test4high|; whose rows equal those of the
    builders of the C code, read from builds of a program that dumps them: in the default build
    $6079$ rows of primes (the rows of $0$ against \verb|sieves0|), $849$ of prime powers and $9455$
    of products, every prime up to $251$ with every residue and every composite modulus up to $255$,
    all their words, the same through the store with the $19\,861$ rows of factors that the products
    built in their tables of rows, and in the build with the chunk $1$ all $16\,383$ rows on their
    first $m$ bit arrays, with a check that one flipped bit or three changed rows show as a
    difference; and whose builds of rows, a line at each set-up of a denominator, at each call of a
    builder by the loop of \verb|sift| and at each row of a factor that a product builds, equal
    those of an instrumented build of the C program line for line on $340$ runs ($17$ curves of
    degrees $3$ to $8$ at heights up to $20\,000$ with five sets of options, all with \verb|-s|,
    four of them with \verb|-k| and one with the reversal, since with the Sturm step on the C code
    skips the set-up of a denominator whose search intervals leave no bits, which belongs to
    stage~H: $1\,003\,520$ set-ups, $114\,535$ rows built, $103\,421$ of primes, $2104$ of prime
    powers and $9010$ of products, and $2788$ rows of factors; the composite moduli from $9$ to $63$
    in the first two stages), with a check that two swapped moduli show as a difference. The first
    stage of either arm agrees for the chunks from $1$ to $12$; at the chunk $17$ the chunked arm
    differs exactly at tails of $16$ bit arrays, which is why the arm is chosen by $\mathrm{chunk}
    \le 16$. Built with other values of the constants (the chunk $1$, $2$, $3$, $5$, $12$ and $17$,
    and $\mathrm{rbaPack}$ $8$, $2$, and $1$ with the chunk $5$), the sketch gives the same rows
    through the store and the same output.
  \item Two reviews of the sketch, each finding checked by a second reader, which found no
    difference from the C code; their findings made the builders store their own rows, as the C code
    does, derived the reference and the kind of builder of a modulus from its factors, kept the
    squares of \verb|examine_power| as words, gave a prime power its table of rows by its prime and
    its exponent, as \verb|pw_idx| keeps its entry, gave line 1800, which sets the builder of a
    prime, back from the sieve entry of a modulus (Definition~\ref{def:makeModulus}) to the table of
    a prime with the exit (Definition~\ref{def:examinePrimeExit}), and added the trace of the builds
    of rows against the C program.
  \item A prototype of all the proofs on the sketch: the whole library builds, standard axioms,
    every proof and every definition at most $20$ lines with no new exception, the statements that
    comparator certifies and \verb|Params.WF| unchanged, the differential tests $3315$ compared with
    none failed and $13$ in \verb|test4high|, and the build with other values of the constants in
    nine variants (the chunk $1$ and $17$ with the other constants changed too, the chunk $2$, $3$,
    $5$ and $12$ alone, $\mathrm{rbaPack}$ $8$ and $2$, and $\mathrm{rbaPack}$ $1$ with the chunk
    $5$); it changed no model file of the sketch.
\end{itemize}
